% ICLR 2027 submission.
%
% The style files in this directory are the OFFICIAL ones, downloaded from
% https://github.com/ICLR/Master-Template/tree/master/iclr2027 -- do not edit
% them. iclr2027_conference_ORIGINAL.tex is the upstream instructions file, kept
% for reference only; it is not part of this submission. The 2027 style differs
% from 2026 only in the running header; the .bst is byte-identical.
%
% Limits (iclr2027_conference_ORIGINAL.tex): a STRICT 9 pages of main text at
% submission, 10 at rebuttal/camera-ready; references, the AI use statement, the
% ethics and reproducibility statements and the appendix do not count.
%
% Every table under tables/ is GENERATED. Do not edit them by hand:
%   python scripts/make_boba_paper_tables.py --analysis output-boba/analysis
% \todo{} marks everything not yet backed by a completed run.
%
%   cd paper && pdflatex main && bibtex main && pdflatex main && pdflatex main

\documentclass{article}
\usepackage{iclr2027_conference,times}

\usepackage{amsmath,amssymb}
\usepackage{booktabs}
\usepackage{graphicx}
\usepackage{placeins}
\usepackage{xcolor}
\usepackage{hyperref}
\usepackage{url}

\newcommand{\todo}[1]{\textcolor{red}{[TODO: #1]}}
\newcommand{\optz}{\mathrm{opt}_z}
\newcommand{\frag}{\mathrm{frag}}

\title{What Noisy Feedback Costs Bayesian\\ Optimization, and the One Number\\ That Predicts It}

% Authors must not appear in the submitted version; keep \iclrfinalcopy
% commented out until camera-ready.
\author{\todo{authors} \\
Affiliation \\
Address \\
\texttt{email}
}

%\iclrfinalcopy % Uncomment for camera-ready ONLY, never for submission.
\begin{document}

\maketitle

\begin{abstract}
Bayesian optimization driven by human feedback is usually evaluated against
regression oracles fitted to archival ratings. That entangles the cost of noisy
feedback with the cost of the oracle being wrong, and a study of one design space
cannot say what generalises. We instead inject four human-plausible error
processes into twenty analytic landscapes with published optima, standardised so
that a unit of error means the same on each: $76{,}800$ paired noisy-versus-clean
runs over four magnitudes, two onsets and twelve acquisition functions. We report
two responses throughout: the optimizer's trajectory, and the design a study
would actually deploy. Error as large as the landscape's own standard deviation,
present from the first rating, costs $14.1\%$ of the available improvement on the
trajectory and $29.7\%$ on the deployed design; arriving at iteration 21 of 50 it
costs $2.8\%$ and $13.6\%$. In trials, one-SD error from the first rating means a median of
$41$ extra trials to match a clean 25-trial study within one percent of the
achievable improvement, and $39\%$ of runs never do within 100. At that magnitude
the landscape and the onset move the cost by $8\times$ and $5\times$, the
acquisition function and the kind of error by $2.8\times$ and $1.3\times$; over
its range the magnitude itself moves it by $23\times$. Measured rater noise in
three human studies falls where the loss is $12$--$20\%$. The expected loss from
a single noisy greedy choice, computable without running an optimizer, explains
nearly as much of the variation as the landscape descriptors and the error
magnitude together, and stays informative beside a descriptor model that lets
magnitude and landscape interact. Decomposing the deployed regret locates that cost: at $1\sigma$, the magnitude at
which the measured rater noise falls, $41.8\%$ $[34, 50]$ of it is
\emph{selection} loss --- the run visited a better design than the one it ships
--- rising to $83.1\%$ $[75, 90]$ when the error is large and arrives late.
Selection loss is identically zero under exact observation, so a noiseless
benchmark never sees it, and no acquisition-side change we tested reduced it.
Consistently,
none of seven acquisition-side changes improves the deployed design, and
augmented expected improvement recovers $41\%$ $[27, 53]$ of the cost on the
search trajectory and none of it on what ships. What works is spending trials on
identification --- a comparative sitting over the last third of the budget gains
$0.051$ $[0.028, 0.072]$, and four confirming trials cut false improvement claims
from $15.7\%$ to $0.5\%$ --- or modelling the fault, worth $22$--$38\%$ at
near-zero price.
Telling the surrogate the true noise variance removes $5\%$ of the excess-weighted
cost and nothing at $1\sigma$ from the first rating ($-7\%$), so the loss is
informational; changing the incumbent definition removes $16\%$ overall but
$12\%$ at $1\sigma$ and up to $41\%$ for one acquisition family, so robustness
rankings rank acquisition-and-incumbent pairs. Every pooled share here is a ratio
of sums that the $5\sigma$ cells dominate, so we give the $1\sigma$ value beside
it throughout. When the person applies the wrong design
rather than the wrong rating, a slip of $15\%$ of each coordinate's range costs
about as much as one-SD rating noise, a record of what was applied removes up
to half of a small slip's cost, and the design the log recommends does
$1.6$--$2.6\times$ worse than the design the person experienced. Of the other process changes we tested, rating early designs twice, replication
and qKG all do worse under error than the standard process, and re-rating the
best designs before deploying helps only a little. On the same grid the fitted-oracle design
reports about a third of the cost, and improving its oracles by up to $0.29$
held-out $R^2$ does not close the gap.
\end{abstract}

\section{Introduction}

Bayesian optimization is the standard tool for optimizing an expensive
black-box function from few evaluations \citep{jones1998efficient,
shahriari2016taking, frazier2018tutorial}, and it is increasingly pointed at
objectives that only a person can evaluate: interface parameters judged for
trust or aesthetics, generated content ranked for quality, robot behaviours
rated for comfort \citep{brochu2010tutorial, gonzalez2017preferential}. The
optimizer assumes each observation is informative about the objective, but a
person is not a sensor. People rate inconsistently, drift, are generous or harsh, make correlated errors,
and sometimes apply the wrong setting altogether. How much that
costs is asked routinely and, we argue, answered with a design that cannot
support the answer.

The standard construction fits a regression model to archival ratings, treats
it as the objective, injects synthetic error into its predictions, and measures
the degradation. This has two problems that no care in the optimization code
can fix. The fitted oracle is itself wrong --- in the motivating studies its
cross-validated $R^2$ as deployed is $0.55$ at best and negative on one
dataset ---
so the measured degradation mixes the injected error with the oracle's
mis-specification, and regret is measured against a random-search estimate of
the optimum that the optimizer can legitimately beat. And the experiment runs on
one design space, so a result such as ``acquisition $X$ is the most robust'' has
no way to say whether it holds anywhere else.

We remove both problems by giving up the human. On an analytic benchmark with a
published optimum the objective is exact, regret is measured from the true
optimum, and the geometry of the landscape can be measured before the
optimization is run and
used to predict what the error will cost. Twenty such landscapes turn ``which acquisition is robust here'' into ``which
property of a problem decides whether feedback error hurts''.

\paragraph{Contributions.}
\begin{enumerate}
  \item A protocol for measuring the cost of feedback error with no fitted
    oracle in the loop: twenty analytic landscapes with verified optima on a
    common scale, four human-plausible error processes, and a model-free control
    whose excess regret is exactly zero.
  \item Measurements of that cost, in regret and in extra trials, and of what
    drives it (Sections~\ref{sec:cost}--\ref{sec:acq}). Magnitude, landscape and
    onset dominate, and a one-shot selection loss computable from the landscape
    and the noise level, without running an optimizer, predicts the cost.
  \item A decomposition of deployed regret into search and selection loss
    (Section~\ref{sec:decomp}): at the magnitude where the measured rater noise
    falls, selection is $41.8\%$ $[34, 50]$ of what error costs, rising to
    $83.1\%$ $[75, 90]$ for a large error arriving late. It is identically zero
    under exact observation, and no acquisition-side change we tested reduced it
    --- which is what reporting regret at the best \emph{observed} point hides.
  \item An account of which remedies can therefore work
    (Section~\ref{sec:budget}, Appendix~\ref{app:budgetneutral}): over
    twenty-four budget-neutral arms, acquisition changes recover nothing on the
    deployed design, while spending a third of the budget on identification, or
    modelling the fault explicitly, recovers $22$--$38\%$, and four confirmation
    trials cut false improvement claims from $15.7\%$ to $0.5\%$.
  \item A matched run of the fitted-oracle design, which reports about a third
    of the cost, a gap that better oracles do not close (Appendix~\ref{sec:fitted}).
  \item An input-error arm, in which the person applies the wrong design rather
    than the wrong rating: a slip of $15\%$ of the range costs about as much as
    one-SD rating noise, and up to half of a small slip's cost is the record
    (Section~\ref{sec:inputerror}). Six further checks, from the surrogate's noise model to a preregistered
    replication on held-out seeds, each change one part of the design
    (Section~\ref{sec:robustness}).
\end{enumerate}

\section{Related work}

\paragraph{Bayesian optimization under noise.}
Noise has been part of the formulation since \citet{jones1998efficient} and is
routine in practice \citep{snoek2012practical}, and several acquisitions are
built for it: augmented expected improvement \citep{huang2006global}, the
knowledge gradient \citep{frazier2009knowledge}, and noisy expected improvement \citep{letham2019constrained, balandat2020botorch}. \citet{picheny2013benchmark} benchmark infill
criteria under noise on analytic problems, the closest antecedent to this work;
we differ in treating the error \emph{process} and its onset as experimental
factors, and in asking which landscape properties predict the outcome rather
than which criterion wins. Heteroscedastic surrogates \citep{kersting2007most} address a case we do not
cover.

\paragraph{Human-in-the-loop and preference-based optimization.}
Where the objective is a person's judgement, the usual response is to change the
observation model: elicit pairwise preferences and fit an ordinal likelihood
\citep{chu2005preference, gonzalez2017preferential}. That handles inconsistent
\emph{scaling} but not inconsistent \emph{ranking}, and it does not tell a
practitioner running a cardinal-rating study what their rater noise will cost.
Appendix~\ref{app:elicitation} measures where the line falls. A comparison is
unharmed by a strictly monotone fault shared across a sitting by construction,
and at equal human cost that buys a great deal against drift and nothing against
a saturating scale, where the map is monotone but not strictly so. We trade the
realism of the objective for
exactness, and map measured rater noise onto the synthetic axis afterwards.

\paragraph{Benchmarks and landscape analysis.}
Results on analytic test functions need not transfer, and suites are small and
idiosyncratic \citep{eggensperger2021hpobench}. Here the twenty landscapes are observations in a regression whose unit is the
landscape, and where that design cannot identify an effect we manipulate the
geometry directly. Predicting algorithm behaviour from
cheap landscape statistics is an old programme in evolutionary computation,
from ruggedness \citep{weinberger1990correlated} and fitness-distance correlation
\citep{jones1995fitness} to their many descendants \citep{malan2013survey}. Our predictor belongs to that family but is defined
relative to the perturbation of interest, a property of the problem and the
error level together.

\section{Setup}

\subsection{Landscapes}
Twenty analytic benchmarks with verified optima, $d = 2$ to $11$: thirteen
classical test functions, six human-performance laws including the
Yerkes--Dodson inverted-U \citep{yerkes1908relation}, and moving peaks
\citep{branke1999memory}; Appendix~\ref{app:descriptors} names them and reports
their geometry. Each objective is affinely
standardised to zero mean and unit variance over its box on a fixed
scrambled-Sobol sample of $2^{16}$ points, so one grid of error magnitudes serves
outputs that span seven orders of magnitude; the map is increasing, so the argmax
is unchanged. One scale is left alone: $\optz$, the number of landscape standard
deviations by which the optimum exceeds the mean (the expected regret of a random
design), runs from $0.76$ to $56.8$, and that $74\times$ spread is what
identifies which scale the cost of error follows (Section~\ref{sec:currency}).

\subsection{Error processes}
\label{sec:errors}
Let $g(x)$ be the standardised objective, $t$ the iteration, $t_0$ the onset and
$\sigma_e$ the swept magnitude, in landscape standard deviations. Observations
are exact for $t \le t_0$; thereafter the optimizer sees $\tilde g_t$:
\begin{align}
  \text{\textsc{gaussian}}\quad & \tilde g_t = g(x_t) + \epsilon_t,
    & \epsilon_t &\sim \mathcal{N}(0, \sigma_e^2) \\
  \text{\textsc{bias}}\quad & \tilde g_t = g(x_t) + \sigma_e + \epsilon_t,
    & \epsilon_t &\sim \mathcal{N}(0, \sigma_e^2) \\
  \text{\textsc{drift}}\quad & \tilde g_t = g(x_t)
    + \sigma_e\tfrac{t - t_0}{T - t_0} + \epsilon_t,
    & \epsilon_t &\sim \mathcal{N}(0, \sigma_e^2) \\
  \text{\textsc{ar}(1)}\quad & \tilde g_t = g(x_t) + e_t,
    \quad e_t = \rho e_{t-1} + \eta_t,
    & \eta_t &\sim \mathcal{N}\!\left(0, \sigma_e^2(1 - \rho^2)\right)
\end{align}
with $\rho = 0.8$ and $e_{t_0+1}$ stationary, so \textsc{ar}(1) has standard
deviation $\sigma_e$ from its first step. \textsc{bias} is a systematically
generous rater, its offset scaled to $\sigma_e$ so that it and \textsc{gaussian}
contrast a systematic with a random error of the same size; \textsc{drift} is
fatigue over a session; \textsc{ar}(1) a rater whose consecutive judgements are
correlated. Noise first touches the observation at $t_0 + 1$, and the first
candidate that can react to it is proposed at $t_0 + 2$. We sweep
$\sigma_e \in \{0.05, 0.25, 1, 5\}$ and $t_0 \in \{0, 20\}$ over $T = 50$
iterations with $5$ initial uniformly random samples, twelve acquisitions and ten
seeds: $79{,}200$ runs, $76{,}800$ of them in noisy/clean pairs.

\subsection{Optimizer and acquisitions}
Every arm uses BoTorch's default single-task GP \citep{balandat2020botorch} ---
RBF kernel with dimension-scaled lengthscale priors, inputs normalised to the
unit cube, outputs standardised, fitted by marginal likelihood. Ten model-based
acquisitions: expected improvement and probability of improvement in analytic,
log \citep{ament2023unexpected} and Monte-Carlo form (EI, LogEI, qEI, PI, LogPI,
qPI), the upper confidence bound in both forms \citep{srinivas2010gaussian}
(UCB, qUCB), noisy expected improvement \citep{letham2019constrained} (qNEI) and
the posterior-mean maximiser (Greedy); and two model-free floors, uniform random
and scrambled Sobol. Improvement-based acquisitions take the maximum posterior
mean at visited points as their incumbent --- the usual choice under noise and,
as Appendix~\ref{sec:incumbent} shows, not a neutral one.

\subsection{Metrics}
\label{sec:metrics}
We report two responses, and they are not interchangeable. \textbf{Search loss}
is the regret of the best design a run has \emph{visited},
$y^\star - \max_{t' \le t} f(x_{t'})$; our primary response is its
\emph{post-onset per-iteration excess}, the noisy run's search loss minus the
identically seeded clean run's, averaged over the iterations the error could have
affected. It is a time average of the optimizer's progress and the response
behind Sections~\ref{sec:cost}--\ref{sec:acq}. \textbf{Deployed loss} is the regret of the design a
study would ship, the one with the best observed rating, at the final trial; it
is search loss plus the selection loss of Section~\ref{sec:decomp}, and about
$2.6\times$ the final search loss. Every table with a search-loss number carries
the deployed number beside it, and a final value is compared only with a final
value. The paired runs share their candidate-proposal random stream, so the
difference isolates the error rather than the seed, and $y^\star$ is a verified
supremum, so no regret is negative. For reporting we divide by $\optz$, the
fraction of the achievable improvement destroyed, and never regress that fraction
on $\optz$ or anything collinear with it, which would manufacture the coefficient
Section~\ref{sec:currency} estimates.

\subsection{Controls}
Three controls, detailed in Appendix~\ref{sec:controls}: the model-free floors
never read an observation, so their excess regret must be identically zero, and
over all $12{,}800$ such comparisons it is; no landscape is ceiling-limited by
its clean run; and a uniform offset is invisible to a GP that standardises its
targets, which \textsc{bias} at onset~0 confirms to within $0.012$ of
\textsc{gaussian} at every magnitude.

\section{How much does feedback error cost?}
\label{sec:cost}

Table~\ref{tab:dose} gives the cost on both responses of
Section~\ref{sec:metrics}. On the optimizer's trajectory, error whose standard
deviation equals the landscape's own, present from the first rating, costs
$14.1\%$ of the improvement that was available, and error five times that large
$22.0\%$. An optimizer that has already located a good region is hard to
dislodge: the same $1\sigma$ error arriving at iteration 21 costs $2.8\%$, a
fifth as much.

\textbf{The design a study would deploy costs far more, and the gap widens with
the error.} At $1\sigma$ from the first rating the shipped design's excess is
$29.7\%$ against the trajectory's $14.1\%$; at $5\sigma$ it is $71.1\%$ against
$22.0\%$, a factor of $3.2$. The late onset, which looks nearly harmless on the
trajectory ($2.8\%$ at $1\sigma$), costs $13.6\%$ on the design that ships. The
trajectory average is bounded by the run's good early iterations; the deployed
value is not, because it is decided at a single moment by a single rating. Every
number below is a trajectory number unless it says otherwise, and
Section~\ref{sec:decomp} takes the deployed value apart.

\begin{table}[h]
\centering
\caption{Fraction of the achievable improvement destroyed by feedback error,
averaged over landscapes and the ten model-based acquisitions. Magnitudes are in
landscape standard deviations. The upper block is the primary response, a time
average over the post-onset trajectory of the best design \emph{visited}; the
lower block is the design the study would \emph{ship}, at the final trial. They
are different quantities and are not comparable across blocks.}
\label{tab:dose}
\begin{tabular}{lrrrr}
\toprule
error present & $0.05\sigma$ & $0.25\sigma$ & $1\sigma$ & $5\sigma$ \\
\midrule
from it.\ 1 & 0.9\% & 5.3\% & 14.1\% & 22.0\% \\
from it.\ 21 & 0.3\% & 1.2\% & 2.8\% & 4.9\% \\
\bottomrule
\end{tabular}

\end{table}

Decomposing the variation by factor, at $1\sigma$ from the first rating unless
stated, puts the practical conclusions in order. Sweeping the magnitude over its
two decades moves the cost by $23\times$; changing the landscape moves it by
$8.1\times$ (Powell $0.038$ to Hartmann-6 $0.303$); moving the onset to
iteration 21, by $5.0\times$; changing the acquisition function, by $2.8\times$
(PI to qUCB); and changing the error \emph{process}, by $1.32\times$ (variance shares over all cells $0.24$, $0.10$, $0.17$, $0.04$, $0.002$: onset
and landscape trade places, the rest stands). The mean cost is $0.148$ for \textsc{gaussian}, $0.152$ for \textsc{drift}, $0.147$ for
\textsc{bias} and $0.115$ for \textsc{ar}(1): \textsc{bias} matches \textsc{gaussian} by
construction at this onset (Appendix~\ref{sec:controls}) and costs $20\%$ more at
the late onset, where the offset arrives as a step, and correlated error is
gentler in every cell, partly absorbed into the surrogate's mean.

\paragraph{Is $1\sigma$ a lot?} We estimate within-rater noise in three
published human studies \todo{cite the three studies in the third person} with
a nearest-neighbour (semivariogram nugget) estimator over repeated ratings by
the same participant, and divide it by $\sigma_f$, the standard deviation of
that study's fitted objective over its design box (Table~\ref{tab:anchor},
Appendix~\ref{app:descriptors}).
Real rating noise lands between $0.74$ and $3.35$ landscape standard
deviations, and participants are noisy from their first rating, so on the
onset-0 row of Table~\ref{tab:dose} these studies would lose $12$--$20\%$ of the
achievable improvement, interpolating log-linearly between columns. The caveat
is that $\sigma_f$ can only be computed through a fitted oracle, and an
under-fit oracle understates $\sigma_f$, so these ratios are if anything high;
this anchors the range of interest rather than measuring it.

\paragraph{How many extra trials?} A practitioner budgets trials, not regret.
For each noisy run we count the trials it needs to come within $1\%$ of the
achievable improvement of its clean twin after $k$ trials, on the budget-100 arm
of Appendix~\ref{sec:designchoices} (\textsc{gaussian} error, six acquisitions;
Table~\ref{tab:extra}; all arms in Appendix~\ref{app:extra}).
Against a clean 25-trial study, error of
$0.05\sigma$ from the first rating costs a median of one extra trial and
$0.25\sigma$ seven; $1\sigma$ costs a median of $41$, a budget of $2.6\times$,
and $39\%$ of runs never get there within 100 trials; at $5\sigma$, $70\%$
never do.  Error from trial 41 costs a 50-trial study no extra trials at the median, though $8$--$12\%$ of runs at $1\sigma$ and above never recover.

\begin{table}[t]
\centering
\caption{Extra trials a noisy run needs to reach the regret its clean twin had
after $k$ trials, within $1\%$ of the achievable improvement (budget-100 arm,
\textsc{gaussian} error, six acquisitions): median over runs, and the share that
never got there within 100 trials.}
\label{tab:extra}
\begin{tabular}{lrr}
\toprule
& error from trial 1, & error from trial 41, \\
error & match a clean 25-trial study & match a clean 50-trial study \\
\midrule
$0.05\sigma$ & 1 (4\%) & 0 (1\%) \\
$0.25\sigma$ & 7 (14\%) & 0 (3\%) \\
$1\sigma$ & 41 (39\%) & 0 (8\%) \\
$5\sigma$ & $>$78 (70\%) & 0 (12\%) \\
\bottomrule
\end{tabular}

\end{table}

\section{Where the cost is: searching or selecting?}
\label{sec:decomp}

The regret reported above is the regret of the design a study would deploy, and
it hides a distinction that decides which remedies can possibly work. Write
$x^\star$ for the optimum, $V$ for the designs a run visited and $d$ for the one
it deploys. Then
\begin{equation}
  \label{eq:decomp}
  \underbrace{f(x^\star) - f(d)}_{\text{deployed regret}}
  \;=\;
  \underbrace{f(x^\star) - \max_{x \in V} f(x)}_{\text{search loss}}
  \;+\;
  \underbrace{\max_{x \in V} f(x) - f(d)}_{\text{selection loss}} .
\end{equation}
The first term is what the optimizer failed to find. The second is what the
ratings failed to identify among what it did find. An acquisition function
addresses the first directly: it chooses where to sample. It can reach the second
only indirectly, by changing which designs are visited and so which set the
ratings must rank --- and none of the seven acquisition-side changes we ran
reduced the deployed loss (Section~\ref{sec:acq}). The selection
term is also exactly zero under exact observation --- the best-rated visited
design is then the best visited design --- which is why it never appears in a
noiseless benchmark and why the field's habit of reporting simple regret at the
best \emph{observed} point is safe there and not here.

We evaluate both terms on all $64{,}000$ noisy runs by re-scoring every visited
design under the loop's own refitted surrogate (Appendix~\ref{app:shiprules}). As
a check rather than an assumption, the largest selection loss over the $2{,}000$
clean runs is $0.00\times10^{0}$, as Equation~\ref{eq:decomp} requires.

\textbf{A large part of what error costs is selection, and how large depends on
the error.} We report the split against the identically seeded clean twin, which
isolates the part the error caused from the part the problem is hard, and we
report it per cell: the share rises steadily with the magnitude and with the
lateness of the onset, from $30.6\%$ $[20, 44]$ at $0.05\sigma$ from the first
rating to $83.1\%$ $[75, 90]$ at $5\sigma$ from trial 21
(Table~\ref{tab:dose}, lower block). \textbf{At $1\sigma$ from the first rating --- the
magnitude at which the measured rater noise of the three human studies falls
(Section~\ref{sec:cost}) --- the excess of $0.297$ splits into $0.173$ of search
loss and $0.124$ of selection loss, so selection is $41.8\%$ $[34, 50]$.} Pooling
every cell gives $61.8\%$ $[53.6, 69.5]$, but that pool is a ratio of sums in
which the $5\sigma$ cells carry $66\%$ of the total and the $1\sigma$ cells
$25\%$, so it is close to the $5\sigma$ answer and we do not use it as the
headline. The pattern is the same for all four error processes.


\textbf{No rule that re-reads the same ratings recovers it.} Four other ship
rules remove essentially none of the selection loss: $+3.9\%$ $[-4, +11]$ for the
posterior mean less two latent SDs, $+0.3\%$ for one, $-1.7\%$ for the mean
rating, $-8.7\%$ for the posterior mean. A likely reason, which we do not test
directly, is that the posterior ranking the visited designs was fitted on the
same corrupted ratings and inherits them. What we can say is that the four rules
differ only in how they re-read data already in hand, and none of them recovers
the term, whereas the procedures of Section~\ref{sec:budget}, which buy new
ratings, do.

\section{What makes an error big?}
\label{sec:currency}

Standardising every objective assumes that a landscape's own standard deviation
is the right unit for its error; the rival is the improvement available, since a
half-SD error is one thing when the optimum stands $0.76$ SDs above a random
design and another at $56.8$. We fit the power law
\begin{equation}
  \mathbb{E}[R] \;=\; A\,\sigma_e^{\beta_c}\,\optz^{\beta_z}
\end{equation}
to the landscape-by-magnitude cell means of the raw excess regret $R$, using quasi-Poisson estimating equations, which keep the cells where noise
helped (Appendix~\ref{app:beneficial}). If the spread is the right unit (\textsc{spread}), $\beta_z = 0$; if the error
relative to the achievable gain is, $R = \optz\, h(\sigma_e/\optz)$, and for a
power-law $h$ that means $\beta_c + \beta_z = 1$ (\textsc{gain}). Intervals are BCa
intervals over resampled landscapes, the unit of replication for a
landscape-level covariate.

\begin{table}[t]
\centering
\caption{Neither scale fits. Exponents of
$\mathbb{E}[R] = A\,\sigma_e^{\beta_c}\optz^{\beta_z}$ fitted to landscape
$\times$ magnitude cell means, with 95\% BCa intervals from 2000 landscape
resamples and a jackknife over landscapes. $\beta_z$ excludes zero in every
condition; $\beta_c + \beta_z - 1$ excludes zero in six of eight.}
\label{tab:currency}
% Eight columns overrun the text width at normal size. Set here rather than in
% tables/currency.tex, which is generated and would lose the fix on a rebuild.
\small
\setlength{\tabcolsep}{4pt}
\begin{tabular}{llrrcrcl}
\toprule
& & & \multicolumn{2}{c}{\textsc{spread}: $\beta_z = 0$} & \multicolumn{2}{c}{\textsc{gain}: $\beta_c + \beta_z = 1$} & \\
\cmidrule(lr){4-5} \cmidrule(lr){6-7}
error process & $t_0$ & $\beta_c$ & $\beta_z$ & 95\% CI & $\beta_c + \beta_z - 1$ & 95\% CI & verdict \\
\midrule
\textsc{gaussian} & 0 & 0.51 & 0.86 & [0.47, 1.13] & +0.37 & [+0.15, +0.70] & neither \\
\textsc{bias} & 0 & 0.51 & 0.85 & [0.49, 1.08] & +0.36 & [+0.17, +0.65] & neither \\
\textsc{drift} & 0 & 0.55 & 0.82 & [0.55, 1.11] & +0.37 & [$-$0.01, +0.63] & gain \\
\textsc{ar}(1) & 0 & 0.66 & 0.85 & [0.60, 1.15] & +0.51 & [+0.14, +0.77] & neither \\
\textsc{gaussian} & 20 & 0.62 & 1.15 & [0.64, 1.48] & +0.76 & [+0.03, +1.03] & neither \\
\textsc{bias} & 20 & 0.70 & 1.14 & [0.61, 1.58] & +0.85 & [+0.04, +1.08] & neither \\
\textsc{drift} & 20 & 0.65 & 1.15 & [0.61, 1.55] & +0.80 & [$-$0.00, +1.14] & gain \\
\textsc{ar}(1) & 20 & 0.71 & 1.14 & [0.64, 1.49] & +0.85 & [+0.13, +1.11] & neither \\
\bottomrule
\end{tabular}

\end{table}

\textsc{spread} fails at both onsets (Table~\ref{tab:currency}): $\beta_z$
excludes zero in all eight conditions, so standardising an objective does not make error magnitudes comparable, and a
noise level reported without the problem's headroom is under-specified. The cost also grows with
$\optz$ faster than \textsc{gain} allows: $\beta_c + \beta_z - 1$ is $+0.36$ to
$+0.51$ from the first observation and $+0.76$ to $+0.85$ from iteration 21,
with intervals excluding zero in six of eight conditions, \textsc{drift}
touching it at both onsets; OLS on the logarithm of the positive cells rejects
\textsc{gain} in all eight. The estimate leans on \texttt{shekel}, the one landscape with $\optz$ above 25;
without it $\beta_z$ is $1.04$--$1.10$ and the gap $0.57$--$0.73$ from the first
observation. Appendix~\ref{sec:manipulation} shows that $\optz$ is not what does the work:
raising a narrow optimum changes nothing, and the dependence runs through how
findable the optimum is.

\section{A one-shot selection loss predicts the cost}
\label{sec:mediation}

Both preceding sections describe the cost of error in terms of properties of the
landscape, but neither gives a practitioner a number. Define $\frag(\sigma_e)$
as the expected loss from a single greedy choice among $m$ quasi-random
candidates when each is observed with $\mathcal{N}(0,\sigma_e^2)$ error:
\begin{equation}
  \frag(\sigma_e) \;=\;
  \mathbb{E}_{\epsilon}\Big[\max_i g_i - g_{\arg\max_i (g_i + \epsilon_i)}\Big],
  \qquad \epsilon_i \sim \mathcal{N}(0,\sigma_e^2),
  \label{eq:frag}
\end{equation}
where $g_i = g(x_i)$ over a fixed Sobol design. It involves no surrogate, no
sequence and no acquisition function; at $m = 256$ it costs about $0.2$ seconds
per landscape.

With onset and error-process indicators in every model, $\frag(\sigma_e)$ explains
$R^2 = 0.52$ of the variation in excess regret across landscape-by-condition cells,
nearly as much as the primary landscape descriptors and the error magnitude together
($0.53$), and it stays informative beside them ($\beta = +0.31$ $[0.18, 0.45]$ in the
interaction model) without making the magnitude redundant. Computed on a GP fitted to
the first ten trials of a clean run rather than on the exact objective, it tracks the
exact value at Spearman $0.94$ and predicts the measured cost as well. Its residuals
are positive in $\optz$ and tail weight: on tall, heavy-tailed landscapes fifty noisy
choices cost more than one predicts (Appendix~\ref{app:mediation}).

\section{Which acquisition function?}
\label{sec:acq}

\begin{table}[t]
\centering
\caption{Acquisition robustness, ranked over landscapes within each condition
and averaged over conditions. ``Absolute loss'' is the deployment number: the
post-onset regret under noise as a percentage of the achievable improvement.}
\label{tab:acq}
\begin{tabular}{lrrr}
\toprule
acquisition & mean rank & excess regret & absolute loss \\
\midrule
qUCB & 3.95 & 4.14\% & 27.9\% \\
qNEI & 4.34 & 4.89\% & 25.2\% \\
UCB & 4.41 & 4.73\% & 27.4\% \\
LogEI & 4.61 & 5.21\% & 26.0\% \\
qEI & 4.73 & 5.37\% & 25.5\% \\
EI & 4.98 & 5.42\% & 26.2\% \\
qPI & 6.29 & 7.90\% & 31.0\% \\
Greedy & 6.56 & 8.11\% & 32.5\% \\
LogPI & 7.48 & 9.19\% & 31.3\% \\
PI & 7.66 & 9.44\% & 31.7\% \\
\midrule
model-free floor & -- & 0.00\% & 41.0\% \\
\bottomrule
\end{tabular}

\end{table}

The confidence-bound family and noisy expected improvement are the most robust,
and the probability-of-improvement family and pure exploitation (Greedy) fill
the bottom four places on both robustness and absolute loss
(Table~\ref{tab:acq}). Averaged over conditions every model-based arm beats the model-free
floor ($41.0\%$); at $5\sigma$ from the first rating the floor matches or beats every
acquisition on this metric, though not on the design it would ship. The Friedman test over the twenty landscapes as blocks reaches
$p < 0.05$ in $28$ of $32$ conditions, and $115$ of $288$ pairwise comparisons against each condition's empirically
best arm survive within-condition Benjamini--Hochberg correction
\citep{benjamini1995controlling}, a comparison that favours the winner.

The landscapes also disagree: the median Kendall's $W$ over the 32 conditions is
$0.247$, and qUCB takes $10$ of $32$ conditions, UCB and qNEI $6$ each, LogEI $5$,
EI $3$, qEI $2$. A study run on one of these landscapes would report a defensible winner that
largely does not transfer, and could not tell. Appendix~\ref{sec:incumbent} shows that part of this ranking is
not about the acquisitions at all.

\textbf{A better search does not ship a better design.} We ran seven
acquisition-side changes as their own arms on the full suite, each keeping the
human trials fixed, scored against the standard process
(Appendix~\ref{app:budgetneutral}). None improves the deployed design: augmented
expected improvement \citep{huang2006global} recovers $+2\%$ $[-6, +11]$ of the
cost of error, a lower-confidence-bound incumbent $-27\%$, Thompson sampling
$-58\%$ $[-81, -40]$, a slip-aware acquisition $-4\%$ to $-35\%$, a
minimum-distance constraint $-13\%$; redistributing the same rating effort across
trials runs from $-7\%$ to $-21\%$. Augmented EI makes the point alone: on the
post-onset trajectory, the response an acquisition is designed to improve, it
recovers $+41\%$ $[+27, +53]$, the best of any arm we ran, and on the design the
study ships it recovers nothing. Under Equation~\ref{eq:decomp} the term it
improves is not the term that dominates.

\section{Buying identification with the trials you have}
\label{sec:budget}

If most of the cost is selection, and no re-reading of the ratings recovers it
(Section~\ref{sec:decomp}), then the only remaining lever is to spend human
trials on identifying rather than searching. We hold the budget at $T=50$ and
spend the last $k$ on one comparative sitting: the top $k$ visited designs by the
ship rule are each shown once, and the best-rated of them is deployed. Everything
before trial $T-k$ is unchanged, so each procedure is replayed exactly from an
existing run log (Appendix~\ref{app:endofstudy}). Error shared across the designs
of one sitting cancels when they are compared, so only the idiosyncratic part
enters the looks; we report both the full idiosyncratic SD and half of it.

\textbf{Identification is worth about a third of the budget.} Against the
standard process, in units of the achievable improvement, a sitting of $k=2$
gains $+0.017$ $[+0.004, +0.029]$, $k=5$ $+0.035$, $k=8$ $+0.042$, $k=12$
$+0.046$, $k=16$ $+0.051$ $[+0.028, +0.072]$, and then falls away: $k=20$
$+0.042$, $k=25$ $+0.022$, $k=30$ $-0.008$. The curve is the inverted U the two
forces imply --- a larger sitting is more likely to contain the best visited
design and leaves fewer trials to find one --- and it turns over at about a third
of $T$. The ceiling, choosing the best $k$ for each run in hindsight, is $+0.145$.
Spending nothing extra and merely changing the ship rule is indistinguishable
from the standard process ($+0.001$, interval $[-0.018, +0.018]$), which is the
selection result of Section~\ref{sec:decomp} in another currency: the trials, not
the rule, are what buys the recovery.

\textbf{Ship the fresh look, not the model.} Of the $19$ variants replayed, the
best takes the top five candidates by the cautious ship rule and deploys
whichever the final sitting rates highest: it recovers $19.7\%$ $[13, 27]$ of the
deployed cost of \textsc{gaussian} error, $33.6\%$ $[25, 41]$ at $5\sigma$ from
trial 21, and $27.9\%$ $[20, 36]$ under an unnoticed slip. Folding the same looks
into the surrogate's posterior rather than trusting them does worse
($11.5\%$), for the reason the selection term exists at all: the posterior was
fitted on the corrupted ratings and carries them forward. It is not universally
safe --- under misclicks it \emph{loses} $14.3\%$ $[-30, -2]$, since a sitting
re-shows the \emph{recorded} design rather than the one that earned the rating,
and below $0.25\sigma$ every variant costs more than the error does.

\textbf{A model-based budget rule inherits the corruption it was meant to fix.}
Choosing $k$ per run from the surrogate's posterior and the run's own rate of
improvement (Appendix~\ref{app:budgetrule}) \emph{loses} to a fixed $k$: its gain
over the standard process is indistinguishable from zero ($+0.004$, interval
$[-0.013, +0.021]$) against the fixed rule's $+0.051$, because in $65\%$ of runs it declines the
sitting for the posterior-mean design, the worst policy of all at $-0.013$. Its
utility is the surrogate's posterior, which was fitted on corrupted ratings, so
it is confident about designs it should not be. Nothing we tried beat fixing $k$
at about a third of the budget.

\section{When the error is in the design, not the rating}
\label{sec:inputerror}

Every error process above corrupts the reported value. People also make the other
mistake: a slider overshoots, the wrong option is pressed. The optimizer proposes
$x$, the person experiences $x' \neq x$, rates it honestly, and the rating is filed
against $x$ --- the right value at the wrong location, which a fitted oracle cannot
score at all, since scoring it needs the objective at $x'$. We sweep a \textsc{slip}
(every trial displaced by $\mathcal{N}(0, s^2)$ per coordinate) and a
\textsc{misclick} (a uniformly random design with probability $p$); two results carry
(Appendix~\ref{app:inputerrorarm}). A slip of $15\%$ of each coordinate's range costs
about what one-SD rating noise costs, so this is not a second-order concern. And
\textbf{the design in the log is not the design that earned its rating}: scored at the
recorded $x$ a practitioner would deploy, the final design's excess regret is $16.4\%$
of the achievable improvement at a $5\%$ slip and $65.8\%$ at $40\%$, against $6.5\%$
and $24.9\%$ for the design the person actually experienced --- a factor of $2.5$ to
$2.6$, of which a record of what was applied removes up to half. The optimizer's
progress survives an input error far better than the record of it, and a deployment
inherits the record.

\section{The conclusions under six further checks}
\label{sec:robustness}

\paragraph{The surrogate's noise estimate is not the problem.} The GP's fitted
noise is an eighth of the injected value for the first ten iterations at
$1\sigma$ and above, yet telling it the true variance removes $5.0\%$ of the
excess-weighted cost and \emph{costs} $7\%$ at $1\sigma$ from the first rating,
with no condition surviving correction (Appendix~\ref{sec:noisediag}).

\paragraph{The incumbent is part of the method.} Switching improvement-based
acquisitions from the posterior-mean incumbent to the observed maximum removes $15.6\%$ of the excess-weighted cost, $12\%$ at $1\sigma$ from the first rating and
$65$--$67\%$ below $0.25\sigma$ ($35$--$41\%$ for the
probability-of-improvement family) but worsens the deployed design: robustness
rankings rank acquisition-and-incumbent pairs (Appendices~\ref{sec:incumbent}
and~\ref{app:adapt}).

\paragraph{$\optz$ is not causal on its own, and dimension is not inert.}
Purpose-built landscapes show that raising a narrow optimum changes nothing, that
a Levy ladder confounds dimension with signal strength, and that with $\optz$
fixed dimension still raises the cost, $1.7\times$ from $d = 7$ to $11$ (Appendix~\ref{sec:manipulation}).

\paragraph{Multi-objective BO is no more fragile.} On six BoTorch problems with published hypervolume optima the dose-response and
onset structures replicate and the four hypervolume acquisitions are
indistinguishable. Early error costs less than on the scalar suite at every
magnitude, but by how much depends on the magnitude and the gap is narrowest
where it matters: a seventh at $0.05\sigma$, a third at $0.25\sigma$ and a half
at $1\sigma$ and $5\sigma$, with the intervals separating only at the two
smallest (Appendix~\ref{sec:multiobjective}).

\paragraph{The budget matters; the rating scale cuts both ways.} With the late
onset held at $40\%$ of the run, the early-to-late cost ratio is $2.3\times$,
$5.1\times$ and $11.7\times$ at $T = 25$, $50$ and $100$; a discretised
rating scale makes small errors costlier and large ones marginally cheaper; and scored against the
standard process, neither qKG nor replication recovers the cost
(Appendices~\ref{sec:designchoices} and~\ref{app:adapt}).

\paragraph{The replication is void; post hoc, the claims hold.} Four claims were preregistered before ten fresh seeds ran; one control failed
(a landscape fell below the headroom threshold), so the run is void by its own
rule, and with that landscape excluded post hoc all four claims clear their
thresholds (Appendix~\ref{sec:confirmatory}).

\paragraph{The fitted-oracle design reports about a third of the cost.} Run on
three archival datasets at matched magnitudes, the data-driven design recovers
$29$--$38\%$ of the cost the analytic arm measures. The gap is not the oracles'
fidelity --- improving them by up to $0.29$ held-out $R^2$ moves the measured
cost by a median of $+0.006$ --- nor the optimum estimate, which cancels in a
difference of two regrets; the three fitted design spaces simply sit at the low
end of the twenty-landscape spread (Appendix~\ref{sec:fitted}).

\section{Discussion}
\label{sec:discussion}

At measured rater-noise levels, feedback error costs a study $12$--$20\%$ of the
improvement it could have had: worth reducing, but not a catastrophe, and not one
extra trials repair, since $39\%$ of runs at $1\sigma$ never match a clean study
with four times its budget. Classifying the kind of error is largely wasted
effort: at $1\sigma$ and above the four processes are within a third of each
other. Estimating its size is not: from a ten-trial pilot and the rater noise,
Equation~\ref{eq:frag} predicts the damage as well as the exact version does.

Most process changes do not help either (Appendix~\ref{app:adapt}): replicating
ratings does worse than the standard process, because the designs it gives up
cost more than the averaging buys, and the observed-max incumbent recovers a
fifth of the trajectory cost while deploying a worse design. Two things remain
unexplained: why the onset effect grows with the
budget (Appendix~\ref{sec:designchoices}), and why the cost grows with $\optz$
faster than \textsc{gain} allows when $\optz$ is not causal
(Appendix~\ref{sec:manipulation}).

\paragraph{Limitations.} The landscapes are analytic, not people: they remove
the surrogate-human confound at the cost of realism, and Appendix~\ref{sec:fitted}
cannot say which design's answer is closer to a person. The rater-noise anchor
runs through a fitted oracle's $\sigma_f$; extra-trial counts are censored at the
budget and depend on the tolerance (Appendix~\ref{app:extra}); the
multi-objective arm is six problems at five seeds, too few to separate the
hypervolume acquisitions; and the preregistered replication is void on a control
failure, so its claims clear their thresholds only after a post-hoc exclusion.

\section*{Reproducibility statement}
The full pipeline is scripted end to end, from the landscape statistics through
the sweep to every table in this paper, and every stochastic component is
explicitly seeded. Regenerating published LogEI and qNEI runs from their recorded settings
reproduces them bit for bit, and the seed order of the model-based acquisitions
has not changed since the sweep. The model-free floors of the main sweep do not
reproduce their recorded noisy observations: their noise seed moved when two
robust baselines were added to the acquisition list after that sweep ran. Their
candidates and regret still reproduce exactly, since a floor never reads an
observation. No reported number
depends on those, and the seed order is now pinned by test. Three further checks are automated. The vendored
benchmark implementations are asserted by test to agree with their source
implementations to $10^{-9}$ relative on random points, so the objective cannot
silently drift. The model-free control of Appendix~\ref{sec:controls} is verified
to be exactly zero in every run. And one recorded optimum in the source
suite, the supremum for a superseded parameter setting, is wrong by $64\%$; we
detected this by searching every box against its recorded value and corrected it
here; the correction and the provenance of the original are
both asserted by test. \todo{repository URL and commit at camera-ready.}

% Required by the ICLR 2027 AI Policy for Authors. The heading, the
% required/not-used/recommended structure and the verification sentence follow
% the template in iclr2027_conference_ORIGINAL.tex. Does not count toward the
% page limit; must not exceed one page.
\subsection*{AI use statement}
In this work, we used generative AI tools (Claude, Anthropic) for the following
tasks with required disclosure: proposing and implementing parts of the
experimental design and analysis; writing and testing most of the simulation and
analysis code; running and monitoring the experiments; and drafting and editing
the manuscript, including the checking of every number in the text against the
analysis outputs. We have not used generative AI tools for
\todo{authors: list the required-disclosure tasks for which no AI was used,
e.g.\ literature search, reviewing, and the choice of research questions}, and
the remaining required-disclosure tasks are not applicable to this work.
Additionally, we used generative AI tools for the recommended-disclosure tasks of
code documentation and copy-editing.

We have reviewed all AI-assisted work. AI-generated code is covered by $697$
automated tests that pin the estimands, the seed order, the negative controls
and hand-computed values, and by cross-checks between independently written
pipelines (for instance, the two computations of the deployed excess agree to
four decimals in every cell, and the table generator refuses to run if they do
not). Three independent audit passes over the code and the manuscript are
recorded in the supplementary material, and the numbers quoted in the text are
verified against the analysis outputs by a script that runs before every build.
The authors set the research questions, chose among the proposed designs, and
take responsibility for the final content of this work, including text, claims
or artifacts produced with the aid of generative AI.

% Recommended by the ICLR 2027 template; does not count toward the page limit.
\subsection*{Ethics statement}
This work collects no new human data. The three human studies it draws on are
published archival rating datasets, used here only to estimate within-rater
noise (Section~\ref{sec:cost}) and as fitted-oracle companions
(Appendix~\ref{sec:fitted}); no participant is identifiable from any quantity we
report. \todo{authors: state the original studies' ethics approval and consent
terms, and confirm that re-analysis is within them.} The analytic landscapes
involve no human participants. We foresee no direct negative societal impact:
the paper's recommendations --- confirm a claimed improvement before reporting
it, and spend part of a human-feedback budget on identification --- reduce the
rate of false improvement claims rather than raise it.

\bibliographystyle{iclr2027_conference}
\bibliography{references}

\clearpage
\appendix
\section{Controls}
\label{sec:controls}
\begin{itemize}
  \item \textbf{Model-free floor.} \texttt{Random} and \texttt{Sobol} never read
    an observation, so their excess regret must be identically zero, and across
    all $12{,}800$ such comparisons the largest absolute value is
    $0.000\times10^{0}$. The analysis raises rather than warns if it is not, and
    the table reads the value from that check rather than printing a constant.
    This negative control shows that no observation noise reaches the candidate
    stream; the floors are then excluded from the robustness ranking, which a
    method that ignores its data would otherwise win.
  \item \textbf{Headroom.} Where the clean run already reaches the optimum a null is
    guaranteed, not earned; the smallest clean-run advantage over the floor is
    $0.176$ (Rosenbrock), so no landscape is ceiling-limited.
  \item \textbf{Uniform offset.} A constant added to every observation is
    invisible to a GP that standardises its targets and uses a posterior-mean
    incumbent. With a shared noise stream the two arms propose identical designs to below
    $10^{-9}$ for the initial design and three model-based candidates (asserted by
    test), after which floating-point differences are amplified; in the sweep, \textsc{bias} and
    \textsc{gaussian} at onset~0 agree to within $0.012$ at every magnitude.
\end{itemize}

\section{The input-error arm}
\label{app:inputerrorarm}

Every error process above corrupts the reported value. People also make the
other mistake: a slider overshoots, the wrong option is pressed, a configuration
is applied to the wrong condition. The optimizer proposes $x$, the person
experiences $x' \neq x$, rates it honestly, and the rating is filed against $x$:
the right value at the wrong location. A fitted oracle cannot score it, since
that needs the objective at $x'$.

A \textsc{slip} displaces every trial by $\delta \sim \mathcal{N}(0, s^2)$ per
coordinate, with $s$ a fraction of that coordinate's range, and clips the result
to the box. A \textsc{misclick} leaves the
trial exact with probability $1 - p$ and otherwise lands on a uniformly random
design. Both magnitudes are swept over $\{0.01, 0.05, 0.15, 0.4\}$ at the same onsets and
budget, with six model-based acquisitions (LogEI, EI, PI, UCB, qUCB, qNEI), the
two floors and five seeds; neither is in landscape standard deviations, so the
comparison with the response-error grid is on the damage axis only.

If the slip goes unnoticed the surrogate trains on $(x, g(x'))$ and the deployed
design is the recorded $x$, so we score it there. A counterfactual arm records
$(x', g(x'))$, as an instrumented interface would, with identical slip draws; the
contrast separates the cost of evaluating the wrong point from that of
mislabelling it. The model-free floors are no longer a zero control: they
evaluate the wrong points without learning, so their excess is the geometric
cost of a slip alone.

\begin{table}[t]
\centering
\caption{Fraction of the achievable improvement destroyed by an input error. A
slip's magnitude is the displacement SD as a fraction of each coordinate's
range; a misclick's is the probability of a random trial. The last three
columns decompose the early-onset slip: \emph{floor}, the model-free arms;
\emph{logged}, a learner whose record holds the point evaluated;
\emph{mislabelling}, the paired difference between the unnoticed slip (second
column) and \emph{logged}. Brackets: 95\% intervals over landscapes.}
\label{tab:inputerror}
\small\setlength{\tabcolsep}{3.5pt}
\begin{tabular}{lrrrrrrr}
\toprule
& \multicolumn{2}{c}{\textsc{slip}} & \multicolumn{2}{c}{\textsc{misclick}} & \multicolumn{3}{c}{\textsc{slip} from it.\ 1, decomposed} \\
\cmidrule(lr){2-3} \cmidrule(lr){4-5} \cmidrule(lr){6-8}
magnitude & \multicolumn{1}{c}{from it.\ 1} & \multicolumn{1}{c}{from it.\ 21} & \multicolumn{1}{c}{from it.\ 1} & \multicolumn{1}{c}{from it.\ 21} & \multicolumn{1}{c}{floor} & \multicolumn{1}{c}{logged} & \multicolumn{1}{c}{mislabelling} \\
\midrule
1\% & 1.5 \makebox[3.0em][l]{\scriptsize [0, 3]} & 1.1 \makebox[3.0em][l]{\scriptsize [0, 2]} & 0.4 \makebox[3.0em][l]{\scriptsize [0, 1]} & 0.1 \makebox[3.0em][l]{\scriptsize [0, 0]} & 0.3 \makebox[3.0em][l]{\scriptsize [0, 1]} & 0.7 \makebox[3.0em][l]{\scriptsize [0, 2]} & 0.8 \makebox[3.0em][l]{\scriptsize [0, 2]} \\
5\% & 5.4 \makebox[3.0em][l]{\scriptsize [3, 8]} & 3.0 \makebox[3.0em][l]{\scriptsize [2, 5]} & 2.8 \makebox[3.0em][l]{\scriptsize [2, 4]} & 0.5 \makebox[3.0em][l]{\scriptsize [0, 1]} & 0.2 \makebox[3.0em][l]{\scriptsize [$-$1, 1]} & 2.8 \makebox[3.0em][l]{\scriptsize [2, 4]} & 2.7 \makebox[3.0em][l]{\scriptsize [1, 4]} \\
15\% & 11.5 \makebox[3.0em][l]{\scriptsize [8, 15]} & 5.0 \makebox[3.0em][l]{\scriptsize [3, 8]} & 6.8 \makebox[3.0em][l]{\scriptsize [5, 9]} & 2.0 \makebox[3.0em][l]{\scriptsize [1, 3]} & 1.2 \makebox[3.0em][l]{\scriptsize [$-$1, 3]} & 7.1 \makebox[3.0em][l]{\scriptsize [5, 10]} & 4.4 \makebox[3.0em][l]{\scriptsize [2, 6]} \\
40\% & 23.2 \makebox[3.0em][l]{\scriptsize [16, 31]} & 6.4 \makebox[3.0em][l]{\scriptsize [4, 10]} & 13.5 \makebox[3.0em][l]{\scriptsize [10, 18]} & 3.5 \makebox[3.0em][l]{\scriptsize [2, 6]} & 6.4 \makebox[3.0em][l]{\scriptsize [0, 13]} & 19.6 \makebox[3.0em][l]{\scriptsize [12, 28]} & 3.6 \makebox[3.0em][l]{\scriptsize [2, 6]} \\
\bottomrule
\end{tabular}

\end{table}

\textbf{A slip costs about as much as a rating error.} A slip of $5\%$ of each
coordinate's range on every trial, from the first, destroys $5.4\%$ of the
achievable improvement; one of $15\%$ destroys $11.5\%$; one of $40\%$, $23.2\%$
(Table~\ref{tab:inputerror}). On this arm's six acquisitions and seeds under \textsc{gaussian} error, the
$0.25\sigma$, $1\sigma$ and $5\sigma$ rating errors of Table~\ref{tab:dose} cost
$3.8\%$, $11.7\%$ and $21.3\%$ (the $T=50$ row of Table~\ref{tab:budget}; pooled:
$5.3\%$, $14.1\%$, $22.0\%$): a $5\%$ slip
costs more than a $0.25\sigma$ rating error, a $15\%$ slip about as much as
$1\sigma$, a $40\%$ slip about as much as $5\sigma$. A misclick is cheaper at
every matched magnitude: even when $40\%$ of trials land somewhere random it
costs $13.5\%$, little more than a $15\%$ slip on every trial and barely half of
a $40\%$ one. A late slip is cheaper than an early one by less than a late rating error:
$1.4$--$3.6\times$ against $5.1\times$ at $1\sigma$ on the same basis.

\textbf{Half of a small slip's cost is the record, not the displacement.}
The decomposition in Table~\ref{tab:inputerror} (in full,
Table~\ref{tab:inputmech}) separates the mechanisms. The model-free floor barely moves, since a random design displaced is still
random; only at $40\%$, where clipping piles evaluations at the box edges, does
it lose $6.4\%$
$[0, 13]$. A learner with a correct record loses far more, $2.8\%$ at a $5\%$ slip and
$19.6\%$ at $40\%$: a displaced proposal is a worse one.
Mislabelling adds $2.7$ points at $5\%$ and $4.4$ at $15\%$, half and two-fifths
of the total (paired intervals $[0.35, 0.61]$ and $[0.26, 0.49]$; not estimable at $1\%$); at $40\%$ the displacement
dominates and the mislabelled share falls to a sixth.

\textbf{The design in the log is not the design that earned its rating.} Scored at the recorded $x$ a practitioner would deploy, the final design's excess
regret is $16.4\%$ of the achievable improvement at a
$5\%$ slip from the first trial, $34.3\%$ at $15\%$ and $65.8\%$ at $40\%$,
against $6.5\%$, $13.4\%$ and $24.9\%$ for the best design the person
experienced (Appendix~\ref{app:inputerror}): a factor of $2.5$ to $2.6$, and
$1.6$ at a $1\%$ slip. The optimizer's progress survives an input error far better than the record of
it, which is what a deployment inherits.

\section{Predicting the cost from one noisy choice}
\label{app:mediation}

The regression behind Section~\ref{sec:mediation}, whose predictor is
Equation~\ref{eq:frag}.

\begin{table}[t]
\centering
\caption{Predicting the cost from one noisy choice. Fitted on landscape
$\times$ condition cell means, predictors $z$-scored, every model also carrying
onset and error-process indicators; the interaction models add
$(\log\sigma_e)^2$ as well. Intervals and $p$-values from a 2000-replicate
bootstrap over landscapes.}
\label{tab:mediation}
% Eight numeric columns with math headers overrun the text width at normal
% size. Set here rather than in tables/mediation.tex, which is generated and
% would lose the fix on the next rebuild.
\small
\setlength{\tabcolsep}{3.2pt}
\begin{tabular}{lrrrrrrrr}
\toprule
model & $\mathrm{frag}(\sigma_e)$ & $\log\sigma_e$ & $\log\mathrm{opt}_z$ & $\log\sigma_e{\times}\log\mathrm{opt}_z$ & $\log$ tail & rugged. & $d$ & $R^2$ \\
\midrule
indicators only & -- & -- & -- & -- & -- & -- & -- & 0.040 \\
$\log\sigma_e$ only & -- & +0.37$^{***}$ & -- & -- & -- & -- & -- & 0.158 \\
$\mathrm{frag}$ only & +0.75$^{***}$ & -- & -- & -- & -- & -- & -- & 0.520 \\
descriptors only & -- & +0.37$^{***}$ & +0.64$^{***}$ & -- & +0.20$^{**}$ & $-$0.09 & $-$0.03 & 0.526 \\
both & +0.66$^{***}$ & $-$0.07 & +0.42$^{**}$ & -- & +0.20$^{**}$ & $-$0.06 & $-$0.03 & 0.701 \\
descriptors + int. & -- & +0.37$^{***}$ & +0.64$^{***}$ & +0.54$^{***}$ & +0.20$^{**}$ & $-$0.09 & $-$0.03 & 0.779 \\
both + int. & +0.31$^{***}$ & +0.17$^{*}$ & +0.54$^{**}$ & +0.40$^{***}$ & +0.20$^{**}$ & $-$0.08 & $-$0.03 & 0.792 \\
\bottomrule
\end{tabular}

\end{table}

With onset and error-process indicators in every model, $\frag(\sigma_e)$
explains $R^2 = 0.52$ of the variation in excess regret across
landscape-by-condition cells ($0.48$ on its own), nearly as much as the primary
landscape descriptors together with the error magnitude ($R^2 = 0.53$;
Table~\ref{tab:mediation}). Jointly they reach $R^2 = 0.70$, and in that additive model the error magnitude
collapses to
$\beta = -0.07$ with $p = 0.44$. That null belongs to the additive specification: with the magnitude-by-$\optz$
interaction Section~\ref{sec:currency} found, and a quadratic in the magnitude,
the descriptor model reaches $R^2 = 0.78$ with no $\frag$ term; $\frag$ still adds to it,
keeping $\beta = +0.31$ $[0.18, 0.45]$ at $R^2 = 0.79$, but the magnitude is
significant again beside it ($-0.38$ $[-0.60, -0.08]$). \textbf{One noisy
choice predicts, on its own, two thirds of what the best descriptor model
explains, and stays informative beside it; it does not make the error magnitude
redundant.} Dimension and ruggedness are not detectably different from zero
here (intervals $[-0.26, 0.12]$ and $[-0.34, 0.15]$), though the fixed-$\optz$
ladder of Appendix~\ref{sec:manipulation} shows dimension does raise the cost.

Equation~\ref{eq:frag} needs the objective at a quasi-random design (here 256
points) and an estimate of rater noise; computed instead on a GP fitted to the first ten trials of a clean run, it tracks
the exact value (Spearman $0.94$) and predicts the measured cost as well
($0.85$ against $0.84$; Appendix~\ref{app:pilot}). Its residuals are positive
in $\optz$ ($+0.42$) and tail weight ($+0.20$): on tall, heavy-tailed landscapes
fifty noisy choices cost more than one predicts.

\section{What the fitted-oracle design would have concluded}
\label{sec:fitted}

The introduction argued that fitting a regression oracle to archival ratings
entangles the cost of feedback error with the cost of the oracle being wrong. To
measure this rather than argue it, we run the data-driven simulator on the three
archival datasets at the same error magnitudes, expressed in each dataset's own
$\sigma_f$ so that ``$1\sigma$'' means the same thing in both arms.

\begin{table}[h]
\centering
\caption{The same experiment, run against exact functions and against oracles
fitted to ratings, under \textsc{gaussian} error. Percentages are post-onset
excess as a fraction of the gap between the optimum and a same-budget model-free
floor (the floor gap), the one normaliser computable in both arms. The fitted
design reports roughly a third of the cost.}
\label{tab:fitted}
\begin{tabular}{lrrrrrrrr}
\toprule
& \multicolumn{4}{c}{error from it.\ 1} & \multicolumn{4}{c}{error from it.\ 21} \\
\cmidrule(lr){2-5} \cmidrule(lr){6-9}
arm & $0.05\sigma$ & $0.25\sigma$ & $1\sigma$ & $5\sigma$ & $0.05\sigma$ & $0.25\sigma$ & $1\sigma$ & $5\sigma$ \\
\midrule
known functions (20) & 10 & 44 & 99 & 148 & 3 & 7 & 10 & 16 \\
fitted oracle (3) & 1 & 11 & 33 & 47 & 1 & 3 & 6 & 8 \\
fitted, no augmentation (3) & 2 & 14 & 38 & 50 & 1 & 4 & 6 & 7 \\
\bottomrule
\end{tabular}

\end{table}

The fitted arm has no verified optimum, so Table~\ref{tab:fitted}'s denominator
uses the best value any arm reached, which is attainable by construction and so
inflates the fitted fraction: the bias works against the result below. Three design
spaces against twenty also widen the intervals.

\textbf{The fitted design understates the cost of feedback error by roughly
threefold.} At $1\sigma$ from the first rating, the known-function arm puts the
loss at $99\%$ of the floor gap (cluster-bootstrap interval $[55, 155]$) and the
fitted arm at $33\%$ ($[5, 75]$): a ratio of $3.0$, interval $[1.1, 17]$. At
$0.05\sigma$ the fitted arm reports $1\%$ ($[-4, 3]$), indistinguishable from
no effect, where the exact functions report $10\%$ ($[4, 18]$). The direction holds in every cell (twofold at the late onset above $0.05\sigma$),
though the intervals overlap at the larger magnitudes.

\textbf{Oracle fidelity does not explain it.} Turning off the deployed jitter
augmentation improves the companion's oracles by up to $0.29$ held-out $R^2$
(\texttt{opticarvis}, $0.52$ to $0.81$). Rerunning the companion arm without it shifts the fraction destroyed by $+0.016$
(median $+0.006$) over the $2{,}400$ shared cells, and the augmentation-free arm
is the worse of the two in $52.6\%$ of them. If the gap were fidelity, it would close as the oracle
improved. Nor can the optimum estimate be the cause: it cancels in the excess, a difference
of two regrets, and enters only the denominator, which biases the other way. The three fitted design spaces sit at the low end of the
twenty-landscape spread (per-landscape fractions at $1\sigma$ run from $0.11$
to $4.8$; the fitted three are $0.05$, $0.19$ and $0.75$), and with three
design spaces the arm cannot say whether that is a property of human
objectives or of these studies.

\section{Ablations}

\subsection{Does the surrogate know it is being lied to?}
\label{sec:noisediag}
The GP fits its observation noise as a free hyperparameter under a prior that
shrinks it, so a measured cost of noise could in principle be the surrogate's
failure to notice the error rather than the information the error destroys. This
is the first thing to rule out, because if it were the explanation then every
number above would be a statement about BoTorch's default prior rather than
about noisy feedback.

The mis-estimation is real, and it lands where it would do the most damage
(Table~\ref{tab:noisediag}). Refitting each surrogate from the sweep's own logs
using only its first $n$ observations, at $1\sigma$ and above the fitted noise standard
deviation is about an eighth of the injected one for the first ten of fifty
iterations (fifteen observations) and has largely recovered by iteration twenty,
while at $0.05\sigma$ it is overestimated by a quarter to a half throughout --- so a diagnostic evaluated only at
$n = 50$ looks reassuring and is misleading. Refitting with a wide noise prior
recovers the injected level far sooner ($0.55$ rather than $0.14$ at $n = 15$,
$1\sigma$), so the shrinkage is the prior's doing and not a limit of the data.

\begin{table}[htbp]
\centering
\caption{Fitted over injected observation-noise standard deviation, as a
function of how many observations the surrogate has seen, for error present from
iteration~1. The estimate is badly wrong exactly while the trajectory is being
decided --- and Table~\ref{tab:knownnoise} shows that this costs almost nothing.}
\label{tab:noisediag}
\begin{tabular}{lrrrr}
\toprule
observations & $0.05\sigma$ & $0.25\sigma$ & $1\sigma$ & $5\sigma$ \\
\midrule
8 & 1.43 & 0.30 & 0.12 & 0.11 \\
15 & 1.50 & 0.41 & 0.14 & 0.10 \\
25 & 1.37 & 0.72 & 0.66 & 0.56 \\
50 & 1.24 & 0.88 & 0.83 & 0.89 \\
\bottomrule
\end{tabular}

\end{table}

It was natural to expect this to explain the onset effect of
Section~\ref{sec:cost}: from the first observation the error arrives in exactly
the window where the noise cannot yet be estimated, whereas at a mid-run onset
the surrogate has a settled length scale and thirty iterations in which to
learn. \textbf{We tested that directly and it is wrong.} Supplying the true
injected variance as \texttt{train\_Yvar} and re-running $4{,}800$ paired cells
removes only $5.0\%$ of the excess-weighted cost ($40\%$ at $0.25\sigma$, not
significant; $3\%$ or less at $1\sigma$ and above), no condition survives FDR
correction
(Table~\ref{tab:knownnoise}), and the onset-0 to onset-20 ratio at $1\sigma$ does
not shrink --- it goes from $4.0$ to $5.2$. Per landscape the largest reduction
is Ackley at $47\%$; per acquisition the effect is inconsistent in sign, helping
PI ($22\%$) and qNEI ($15\%$) but \emph{hurting} qUCB ($-29\%$).

Two conclusions follow. Bayesian optimization is remarkably insensitive to its
own noise hyperparameter in this regime, despite getting it badly wrong for the
first fifth of the budget --- worth knowing on its own, and a caution
against reading too much into noise-estimation diagnostics. And the cost this
paper measures is the information cost of the error, not an artefact of a
mis-specified surrogate, which is what licenses reading Table~\ref{tab:dose} as
a property of the problem. The onset effect itself is therefore left
unexplained; we return to it in Section~\ref{sec:discussion}.

\begin{table}[htbp]
\centering
\caption{Telling the surrogate the truth about its noise. Paired within
(landscape, acquisition, magnitude, onset, seed); positive $\Delta$ means the
treatment suffers more. Nothing survives FDR correction.}
\label{tab:knownnoise}
\begin{tabular}{rrrrrrrr}
\toprule
$\sigma_e$ & $t_0$ & reference & treatment & $\Delta$ & share & $d_z$ & $p_{\text{FDR}}$ \\
\midrule
0.05 & 0 & 0.065 & 0.077 & +0.012 & $-$19\% & 0.03 & 0.888 \\
0.25 & 0 & 0.271 & 0.166 & $-$0.105 & +39\% & $-$0.18 & 0.888 \\
1 & 0 & 0.790 & 0.848 & +0.058 & $-$7\% & 0.10 & 0.888 \\
5 & 0 & 1.484 & 1.463 & $-$0.021 & +1\% & $-$0.04 & 0.888 \\
0.05 & 20 & 0.002 & 0.001 & $-$0.002 & +71\% & $-$0.02 & 0.913 \\
0.25 & 20 & 0.072 & 0.041 & $-$0.031 & +43\% & $-$0.20 & 0.888 \\
1 & 20 & 0.197 & 0.164 & $-$0.034 & +17\% & $-$0.19 & 0.888 \\
5 & 20 & 0.531 & 0.483 & $-$0.048 & +9\% & $-$0.16 & 0.888 \\
\bottomrule
\end{tabular}

\end{table}

\FloatBarrier
\subsection{Is it the acquisition or the incumbent?}
\label{sec:incumbent}
$\mathrm{best}_f$ is the maximum posterior mean at visited points, which shrinks
under noise and so makes improvement-based acquisitions automatically more
exploratory. The acquisitions differ in their exposure to this: LogEI, EI, PI,
LogPI and qEI consume $\mathrm{best}_f$; UCB does not; qNEI ignores it entirely.
A ranking of acquisitions by robustness is therefore partly a ranking of their
sensitivity to a design choice that is not part of the acquisition at all.
Re-running with the observed maximum as the incumbent removes $15.6\%$ of the
excess-weighted cost overall ($14\%$ at $1\sigma$, two thirds or more at
$0.25\sigma$ and below), with four of eight conditions significant after FDR
correction (Table~\ref{tab:incumbent}).

The per-acquisition breakdown is in Table~\ref{tab:incumbentacq}. UCB
and qNEI move by \emph{exactly} zero, as they must --- a control that validates
the pairing at machine precision. The probability-of-improvement family moves
most: LogPI loses $41\%$ of its measured cost and PI $35\%$, against $17\%$ for
EI, $2\%$ for LogEI and $-2\%$ for qEI. So a substantial part of ``PI is fragile
under noise'' is really ``the posterior-mean incumbent is fragile under noise,
and PI is the acquisition most exposed to it''. Every acquisition-robustness
ranking here, including Table~\ref{tab:acq}, is a ranking of
acquisition-and-incumbent pairs, and the same is true of the applied results it
is modelled on. Reporting the incumbent definition alongside the acquisition should be standard.
Scored against the standard process's clean run rather than its own, the switch
still recovers a fifth of the trajectory cost, but the design it deploys is
worse (Appendix~\ref{app:adapt}).

\begin{table}[htbp]
\centering
\caption{Observed-maximum against posterior-mean incumbent, same pairing as
Table~\ref{tab:knownnoise}.}
\label{tab:incumbent}
\begin{tabular}{rrrrrrrr}
\toprule
$\sigma_e$ & $t_0$ & reference & treatment & $\Delta$ & share & $d_z$ & $p_{\text{FDR}}$ \\
\midrule
0.05 & 0 & 0.010 & 0.003 & $-$0.007 & +72\% & $-$0.56 & 0.019 \\
0.25 & 0 & 0.050 & 0.026 & $-$0.024 & +48\% & $-$1.39 & $<$0.001 \\
1 & 0 & 0.132 & 0.100 & $-$0.032 & +24\% & $-$0.93 & 0.004 \\
5 & 0 & 0.223 & 0.200 & $-$0.022 & +10\% & $-$0.71 & 0.011 \\
0.05 & 20 & 0.003 & 0.001 & $-$0.002 & +81\% & $-$0.24 & 0.142 \\
0.25 & 20 & 0.014 & 0.005 & $-$0.009 & +64\% & $-$0.88 & 0.002 \\
1 & 20 & 0.027 & 0.021 & $-$0.005 & +20\% & $-$0.50 & 0.032 \\
5 & 20 & 0.049 & 0.048 & $-$0.000 & +0\% & $-$0.01 & 0.421 \\
\bottomrule
\end{tabular}

\end{table}

\begin{table}[htbp]
\centering
\caption{The incumbent arm by acquisition. UCB and qNEI never read
$\mathrm{best}_f$ and move by exactly zero, which validates the pairing; the
probability-of-improvement family moves most.}
\label{tab:incumbentacq}
\begin{tabular}{lrrrr}
\toprule
acquisition & reference & treatment & $\Delta$ & share removed \\
\midrule
LogPI & 0.555 & 0.328 & $-$0.227 & +41\% \\
PI & 0.572 & 0.373 & $-$0.199 & +35\% \\
EI & 0.459 & 0.382 & $-$0.077 & +17\% \\
LogEI & 0.493 & 0.481 & $-$0.011 & +2\% \\
qNEI & 0.378 & 0.378 & +0.000 & +0\% \\
UCB & 0.346 & 0.346 & +0.000 & +0\% \\
qEI & 0.445 & 0.454 & +0.009 & $-$2\% \\
\bottomrule
\end{tabular}

\end{table}

\FloatBarrier
\subsection{Breaking the descriptor confound}
\label{sec:manipulation}
In our suite $\optz$, sparsity and skew correlate at about $0.9$, and for a
bounded function with an isolated optimum that is close to a structural
identity, so Section~\ref{sec:mediation}'s descriptor terms cannot be
separated observationally, however many real benchmarks are added. Two
purpose-built families vary them by construction: a cos-field background plus one
spike, whose amplitude moves $\optz$ and whose width moves sparsity; and the Levy
function at $d \in \{4, 7, 11\}$, one landscape family across a dimension ladder
(Table~\ref{tab:manipulation}).

\textbf{$\optz$ on its own is not causal.} Holding the spike narrow and raising
its amplitude by $4\times$ moves $\optz$ from $5.0$ to $21.7$ and changes the
excess regret at $1\sigma$ from $0.246$ to $0.245$ --- no effect at all. Widen
the spike and the same amplitude change raises excess from $0.495$ to $1.888$.
What the observational $\beta_{\optz} = +0.42$ picks up is therefore an
interaction, not a main effect: a taller optimum costs more to miss only when it
is broad enough that the optimizer would otherwise have found it. On a true
needle, the optimizer was never going to find the optimum with or without noise,
so there is nothing for the noise to destroy. It also means that the $\optz$ exponent of Section~\ref{sec:currency}
describes the suite rather than a mechanism: it is not $\optz$ as such that
makes an error expensive.

\textbf{Dimension, and a within-family ladder that corrects itself.} Within the
Levy family, excess regret at $1\sigma$ rises monotonically from $0.114$ at
$d = 4$ to $0.222$ at $d = 7$ and $0.264$ at $d = 11$, while the cross-suite
regression puts $\beta_d$ at $-0.03$ and not significant. Read on its own that
looks like a real effect the observational analysis missed. It is not, or not
mostly: $\optz$ also rises along the Levy ladder, $1.53$ to $2.07$ to $2.62$, so
that comparison varies dimension and signal strength together.

The \texttt{bump} ladder separates them. Amplitude and width are chosen at each
dimension so that $\optz$ is $9.0$ and the spike occupies a $10^{-6}$ fraction of
the box at all three, leaving dimension as the only difference. In the table's
units, excess regret is $0.28$ at $d = 4$, $0.25$ at $d = 7$ and $0.43$ at
$d = 11$. The $d = 4$ rung has headroom $0.098$, just below the $0.10$ screen
applied everywhere else, so the comparison that clears the screen is $d = 7$
to $11$: there the fixed-$\optz$ ladder rises $1.7\times$ where Levy rises
$1.2\times$. \textbf{Dimension does raise the cost when $\optz$ is held fixed;
what the Levy ladder cannot say is how much of its $2.3\times$ is dimension
rather than signal strength.} What rises monotonically with dimension on both
ladders is the inference metric, the true value of the design a practitioner
would deploy: $0.52$, $0.70$, $0.99$ on the fixed-$\optz$ ladder. Higher
dimension makes the progress harder to identify from noisy observations at the
end as well as harder to make.

The methodological lesson is that a within-family ladder is not automatically a
manipulation of the property it is indexed by: the Levy ladder varied signal
strength along with dimension, and ours became a manipulation of dimension only
when it was built to hold $\optz$ fixed.

\begin{table}[htbp]
\centering
\caption{Purpose-built landscapes. Raw post-onset excess regret in landscape
standard deviations, \textsc{gaussian} error from iteration~1. Amplitude and width are varied independently in the top block; the middle block
is one function family at three dimensions; the bottom block holds $\optz = 9.0$
and the spike volume fixed across three dimensions.}
\label{tab:manipulation}
\begin{tabular}{lrrrrrrr}
\toprule
landscape & $d$ & $\mathrm{opt}_z$ & sparsity & $0.05\sigma$ & $0.25\sigma$ & $1\sigma$ & $5\sigma$ \\
\midrule
\texttt{bump\_a4\_w0.05} & 4 & 5.0 & $<$1.5e-05 & 0.007 & 0.152 & 0.246 & 0.324 \\
\texttt{bump\_a16\_w0.05} & 4 & 21.7 & $<$1.5e-05 & 0.014 & 0.026 & 0.245 & 0.262 \\
\addlinespace
\texttt{bump\_a4\_w0.15} & 4 & 5.3 & 3.8e-04 & 0.133 & 0.289 & 0.495 & 0.704 \\
\texttt{bump\_a16\_w0.15} & 4 & 12.3 & 1.8e-04 & 0.298 & 0.730 & 1.888 & 3.173 \\
\midrule
\texttt{levy\_4d} & 4 & 1.5 & 2.6e-01 & 0.006 & 0.048 & 0.114 & 0.200 \\
\texttt{levy\_7d} & 7 & 2.1 & 7.3e-02 & 0.051 & 0.115 & 0.222 & 0.305 \\
\texttt{levy\_10} & 11 & 2.6 & 1.3e-02 & 0.001 & 0.147 & 0.264 & 0.481 \\
\midrule
\texttt{bump\_d4} & 4 & 9.0 & $<$1.5e-05 & 0.001 & 0.121 & 0.279 & 0.327 \\
\texttt{bump\_d7} & 7 & 9.0 & $<$1.5e-05 & $-$0.001 & 0.103 & 0.248 & 0.428 \\
\texttt{bump\_d11} & 11 & 9.0 & $<$1.5e-05 & 0.142 & 0.215 & 0.427 & 0.734 \\
\bottomrule
\end{tabular}

\end{table}

\FloatBarrier
\section{Does any of this hold for multiple objectives?}
\label{sec:multiobjective}

The applied setting is frequently multi-objective \citep{daulton2021parallel}
and everything above is scalar, so the generalisation cannot be assumed. We
repeat the design on seven BoTorch problems with a published maximum
hypervolume --- BraninCurrin, ZDT1, DTLZ2, DH2, VehicleSafety, Penicillin and
CarSideImpact --- spanning two to four objectives and two to seven dimensions,
with the four hypervolume acquisitions (qEHVI, qNEHVI and their log variants) and
the same two model-free floors, under \textsc{gaussian} error, at five seeds: $1{,}680$ paired runs, of
which $960$ are model-based noisy-versus-clean comparisons on the problems that
survive the headroom screen below.

Standardisation transfers exactly rather than approximately. Each objective is
mapped to zero mean and unit variance over a fixed scrambled-Sobol sample of
$2^{15}$ points, drawn as for the scalar suite, and under a per-objective affine map the hypervolume scales by
$\prod_m \sigma_m$, so the standardised maximum follows in closed form from the
published one. We verify this rather than assert it: across the seven problems
the largest relative error between the closed-form standardised maximum and a
direct recomputation is $4.7 \times 10^{-16}$.

\paragraph{Controls.} The model-free floor is again identically zero across all
$560$ such comparisons, so no observation noise leaks into the candidate stream
in the hypervolume code path either. No run beats its published maximum
hypervolume. Acquisition-optimization fallbacks total zero, which we report
rather than assume: a hypervolume acquisition that fails to construct falls back
to random sampling, and because that produces a plausible-looking curve rather
than an error, the arm has to be read off the fallback counter and not off the
plots.

\paragraph{One problem is dropped, by the same screen as before.} DH2's
reference point leaves a $32{,}768$-point Sobol sample with a hypervolume of
exactly zero, and its best learner does not improve on the model-free floor at
all. There is nothing for error to destroy, so its null is guaranteed by
construction; it is excluded. The remaining six pass, with Penicillin weakest at
$0.197$ --- comparable to Rosenbrock at $0.176$, the weakest admissible scalar
landscape.

\begin{table}[htbp]
\centering
\caption{The two arms on a common footing. Both columns are post-onset excess
divided by the gap between the optimum and a same-budget model-free floor, as a
percentage; brackets are 95\% cluster-bootstrap intervals over
landscapes/problems. This is \emph{not} the normalisation used in
Table~\ref{tab:dose}, and the two must not be read against each other; the scalar column is recomputed from the same runs as the known-function row of
Table~\ref{tab:fitted}, under \textsc{gaussian} error.}
\label{tab:mo}
\begin{tabular}{lrrrr}
\toprule
& \multicolumn{2}{c}{scalar (20 landscapes)} & \multicolumn{2}{c}{multi-objective (6 problems)} \\
\cmidrule(lr){2-3} \cmidrule(lr){4-5}
error & \multicolumn{1}{c}{from trial 1} & \multicolumn{1}{c}{from trial 21} & \multicolumn{1}{c}{from trial 1} & \multicolumn{1}{c}{from trial 21} \\
\midrule
$0.05\sigma$ & 10.2 \makebox[3.4em][l]{\scriptsize [4, 19]} & 2.8 \makebox[3.4em][l]{\scriptsize [1, 5]} & 1.4 \makebox[3.4em][l]{\scriptsize [0, 3]} & 0.5 \makebox[3.4em][l]{\scriptsize [$-$1, 2]} \\
$0.25\sigma$ & 44.3 \makebox[3.4em][l]{\scriptsize [22, 71]} & 6.5 \makebox[3.4em][l]{\scriptsize [4, 10]} & 13.6 \makebox[3.4em][l]{\scriptsize [8, 19]} & 3.7 \makebox[3.4em][l]{\scriptsize [0, 10]} \\
$1\sigma$ & 99.0 \makebox[3.4em][l]{\scriptsize [54, 154]} & 10.3 \makebox[3.4em][l]{\scriptsize [7, 13]} & 49.1 \makebox[3.4em][l]{\scriptsize [42, 56]} & 10.7 \makebox[3.4em][l]{\scriptsize [2, 24]} \\
$5\sigma$ & 148.4 \makebox[3.4em][l]{\scriptsize [76, 239]} & 16.5 \makebox[3.4em][l]{\scriptsize [11, 22]} & 72.3 \makebox[3.4em][l]{\scriptsize [55, 91]} & 13.0 \makebox[3.4em][l]{\scriptsize [3, 27]} \\
\bottomrule
\end{tabular}

\end{table}

\paragraph{The comparison needs one piece of care.} Elsewhere the paper divides
excess regret by $\optz$, a z-score of the optimum against a random design's
mean. A Pareto problem has no scalar optimum and therefore no $\optz$, and the
natural hypervolume analogue --- the gap between the published maximum and what a
model-free design of the same budget reaches --- is a numerically much smaller
denominator. Quoting one arm's published fraction against the other's would
manufacture a threefold difference out of the choice of divisor alone. Both
columns of Table~\ref{tab:mo} are therefore recomputed on the floor denominator,
identically; and we also report the onset ratio, which is a ratio of two costs on
the same problem and so does not depend on the denominator at all.

\textbf{The two headline structures replicate.} Cost rises monotonically with
error magnitude and falls sharply when the error arrives late, on both arms. At
$1\sigma$ from the first rating the multi-objective arm loses $49.1\%$ of the hypervolume floor gap $[42, 56]$ against the scalar arm's
$99.0\%$ $[54, 154]$; arriving at iteration~21 the two are indistinguishable,
$10.7\%$ $[2, 24]$ against $10.3\%$ $[7, 13]$.

\textbf{There is no evidence multi-objective BO is the more fragile of the two,
and some that it is less.} The point estimates for early error are half the scalar arm's or less at every
magnitude, a seventh at $0.05\sigma$, a third at $0.25\sigma$ and half at
$1\sigma$ and $5\sigma$, and the intervals separate at $0.05\sigma$
and $0.25\sigma$; at $1\sigma$ they overlap slightly and at $5\sigma$
substantially, so this is a consistent direction rather than a
clean separation at the top of the range. The onset ratio is correspondingly
flatter: $3.0$--$5.5\times$ across magnitudes against the scalar arm's
$3.6$--$9.6\times$. A mechanism consistent with both observations, and with
Appendix~\ref{sec:incumbent}, is that a hypervolume objective has no single
incumbent to corrupt, so the \emph{search} keeps finding good designs. The front
the study would \emph{report} is another matter: at $1\sigma$ from the first
rating about half of it is truly dominated ($5.2$ of $10.9$ designs), it misses
most of the truly non-dominated designs evaluated, and its hypervolume shortfall
exceeds that of the evaluated set in every cell.

\begin{table}[htbp]
\centering
\caption{Hypervolume acquisitions, ranked over problems within each condition.
Absolute loss is the fraction of the hypervolume floor gap lost under
noise. The four are not distinguishable.}
\label{tab:moacq}
\begin{tabular}{lrr}
\toprule
acquisition & mean rank & absolute loss \\
\midrule
qNEHVI & 2.40 & 20.2\% \\
qLogEHVI & 2.44 & 20.4\% \\
qEHVI & 2.56 & 20.7\% \\
qLogNEHVI & 2.60 & 20.8\% \\
\bottomrule
\end{tabular}

\end{table}

\textbf{The choice among hypervolume acquisitions does not appear to matter.}
Mean ranks span $2.40$ to $2.60$ on a scale where $2.5$ is a tie, absolute losses
span $20.2\%$ to $20.8\%$, and the Friedman test reaches $p < 0.05$ in $0$ of
$8$ conditions (Table~\ref{tab:moacq}). Kendall's $W$ over
condition-averaged losses is $0.056$, against $0.52$ for the scalar suite computed
the same way, and each of the four acquisitions is the most
robust on at least one of the six problems. Six problems at five seeds is far
less power than twenty at ten, so this is ``no detectable
difference'' rather than ``no difference''. It does, though, point the same way as
Section~\ref{sec:acq}: the disagreement between design spaces about which
acquisition is most robust is, if anything, sharper here.

\paragraph{What does not transfer.} Equation~\ref{eq:frag} is defined for a
scalar objective: it is the loss from one greedy pick among $m$ candidates, and
on a Pareto problem there is no single greedy pick whose loss it measures. Section~\ref{sec:mediation}'s result is therefore untested in this arm rather than confirmed or
refuted by it. A hypervolume analogue would have to define the one-shot loss over
hypervolume contributions and take a stand on how per-objective error propagates
into them; that is a paper of its own, and we do not claim
Table~\ref{tab:mediation} extends until someone writes it.

\FloatBarrier
\section{Three design choices that could have produced the result}
\label{sec:designchoices}

Every number above was measured at a fixed budget, against an unbounded
continuous response, using standard acquisition functions. Each of those is a
choice, and each could be doing the work. The three arms here ran
\textsc{gaussian} error on seeds 7--11 with six acquisitions (LogEI, EI, PI, UCB,
qUCB, qNEI) or, for the robust baselines, the two named; every reference row is
the main sweep restricted to the same error process, acquisitions and seeds.

\begin{table}[htbp]
\centering
\caption{Budget, with the late onset held at $40\%$ of the run in each case and
the reference row restricted to the arm's own error process, acquisitions and
seeds, so that budget is the only thing that changes. Percentages are the
fraction of the achievable improvement destroyed.}
\label{tab:budget}
\begin{tabular}{llrrrr}
\toprule
budget & error present & $0.05\sigma$ & $0.25\sigma$ & $1\sigma$ & $5\sigma$ \\
\midrule
$T=25$ & from it.\ 1 & 0.6 & 3.2 & 9.4 & 15.4 \\
 & from it.\ 11 & 0.3 & 1.3 & 4.1 & 8.0 \\
 & ratio & 2.2$\times$ & 2.4$\times$ & 2.3$\times$ & 1.9$\times$ \\
\midrule
$T=50$ & from it.\ 1 & 0.6 & 3.8 & 11.7 & 21.3 \\
 & from it.\ 21 & 0.1 & 1.0 & 2.3 & 4.7 \\
 & ratio & 4.4$\times$ & 3.9$\times$ & 5.1$\times$ & 4.6$\times$ \\
\midrule
$T=100$ & from it.\ 1 & 0.2 & 3.4 & 11.1 & 23.9 \\
 & from it.\ 41 & 0.1 & 0.4 & 0.9 & 2.2 \\
 & ratio & 1.4$\times$ & 9.2$\times$ & 11.7$\times$ & 10.8$\times$ \\
\bottomrule
\end{tabular}

\end{table}

\paragraph{Budget, and the onset effect is not budget-free.} Re-running at $T=25$
and $T=100$ with the onset held at the same fraction of the run gives
Table~\ref{tab:budget}. Early-error cost is roughly flat in budget --- $9.4\%$,
$11.7\%$, $11.1\%$ at $1\sigma$ --- but late-error cost falls steeply, from
$4.1\%$ at $T=25$ to $2.3\%$ at $T=50$ to $0.9\%$ at $T=100$. The onset ratio
therefore grows with the budget, from $2.3\times$ to $5.1\times$ to $11.7\times$.
This is the clearest limitation on the headline: \textbf{the $5\times$ onset
effect is a statement about a fifty-iteration study.} The direction is
unsurprising, since more remaining iterations means more chance to recover, but
the size is not: a fourfold change in budget moves the ratio by five. We looked
for a rescaling that would let one budget speak for another and did not find
one. Late-onset cost per remaining iteration falls by about $17\times$ across the
three budgets, and late-onset cost times remaining iterations is non-monotone
($0.62$, $0.69$, $0.57$), so neither is the invariant. The early-error number is
the transferable one; a study at a different budget has to measure its own onset
effect.

\begin{table}[htbp]
\centering
\caption{A bounded, discrete response. The instrument arm clips to $\pm 8\sigma$
and rounds to $0.55\sigma$, the range and step implied by the three archival
studies. Both rows: \textsc{gaussian} error, six acquisitions, five seeds.}
\label{tab:instrument}
\begin{tabular}{llrrrr}
\toprule
response scale & error present & $0.05\sigma$ & $0.25\sigma$ & $1\sigma$ & $5\sigma$ \\
\midrule
unbounded, continuous & from it.\ 1 & 0.6 & 3.8 & 11.7 & 21.3 \\
 & from it.\ 21 & 0.1 & 1.0 & 2.3 & 4.7 \\
bounded, discrete & from it.\ 1 & 3.8 & 5.0 & 11.7 & 20.8 \\
 & from it.\ 21 & 1.3 & 1.4 & 2.5 & 4.6 \\
\bottomrule
\end{tabular}

\end{table}

\paragraph{A real rating scale, which cuts both ways.} The synthetic objective is
unbounded and continuous; a person returns a bounded, discrete rating. The three
archival instruments have steps of $0.55$, $0.54$ and $1.12\sigma$ and span
sixteen to twenty $\sigma$, so the arm clips to $\pm 8\sigma$ and rounds to
$0.55\sigma$ --- the two five-item instruments, in the study's own currency.
Table~\ref{tab:instrument} shows the two effects separately. At $0.05\sigma$ the
discretised arm is six times worse ($3.8\%$ against $0.6\%$), because rounding
to a $0.55\sigma$ grid is itself an error of about $0.16\sigma$ and swamps an
injected error far smaller than the step; at $0.25\sigma$ it is still worse
($5.0\%$ against $3.8\%$); at $1\sigma$ the two are indistinguishable
($11.7\%$ against $11.7\%$); and at $5\sigma$ it is slightly better ($20.8\%$
against $21.3\%$), because clipping bounds how far one bad rating can move the
incumbent. A rating scale is a noise floor at the bottom of the range and, at
the top, a marginal safety rail.

\begin{table}[htbp]
\centering
\caption{Robust baselines. qKG is the knowledge gradient; replication spends
every second evaluation re-asking about the incumbent. Compared at a fixed number
of evaluations, so replication buys averaging with half the designs. The
standard row is the ten model-based acquisitions of the main sweep on the same
error process and seeds.}
\label{tab:robust}
\begin{tabular}{llrrrr}
\toprule
acquisitions & error present & $0.05\sigma$ & $0.25\sigma$ & $1\sigma$ & $5\sigma$ \\
\midrule
ten standard & from it.\ 1 & 1.1 & 5.6 & 13.9 & 21.5 \\
 & from it.\ 21 & 0.3 & 1.3 & 2.7 & 4.5 \\
qKG and replication & from it.\ 1 & 0.3 & 2.2 & 9.2 & 16.2 \\
 & from it.\ 21 & 0.2 & 0.1 & 2.2 & 4.8 \\
\bottomrule
\end{tabular}

\end{table}

\paragraph{Methods built for noise do not recover the cost.} All ten
model-based acquisitions above are ordinary ones run under noise. Two baselines
built for it: the knowledge gradient, which values the improvement in the
posterior \emph{maximum} rather than in any observed value, and spending every
second trial re-asking about the incumbent. Scored the way the rest of this
appendix is, against each method's own clean run, they lose less: $9.2\%$ of the
achievable improvement at $1\sigma$ from the first rating against the standard
acquisitions' $13.9\%$, and $16.2\%$ against $21.5\%$ at $5\sigma$
(Table~\ref{tab:robust}). That comparison flatters them. Replication re-rates the
incumbent in its clean run too, which carries no information, and qKG's clean
run is itself worse than the standard acquisitions', so part of the smaller
excess is a worse reference. Scored against the standard process's clean run
(Appendix~\ref{app:adapt}), replication recovers none of the cost ($-61\%$ on
the trajectory, $-5\%$ on the deployed design) and qKG recovers some of it at the
early onset ($+12\%$ to $+43\%$ at $0.25\sigma$ and above) but more than loses it
at the late onset and in the design it recommends ($-32\%$). What remains is the
information the error destroys, which is what Equation~\ref{eq:frag} measures.

\FloatBarrier
\section{Held-out seeds: a void replication and its sensitivity analysis}
\label{sec:confirmatory}

Everything above was selected and estimated on the same ten seeds. Before
running any fresh ones we wrote down four claims, the statistics that would test
them, the thresholds each had to clear, and four control conditions that would
void the run, then ran ten fresh seeds under \textsc{gaussian} error only, as
preregistered.

\begin{table}[htbp]
\centering
\caption{Preregistered quantities on the screening seeds and on ten fresh ones.
Both columns use \textsc{gaussian} error only and exclude Rosenbrock, so the
screening column differs slightly from the pooled twenty-landscape values of
Sections~\ref{sec:cost}--\ref{sec:acq}. The replication line reports the
\emph{sensitivity} analysis described in the text, not the preregistered run,
which is void.}
\label{tab:confirmatory}
\begin{tabular}{lrr}
\toprule
preregistered quantity & screening & fresh seeds \\
\midrule
H1 $\beta(\frag)$ & +0.65 & +0.69 \\
H1 $\beta(\log_{10}\sigma_e)$ & $-$0.06 & $-$0.07 \\
H2 onset ratio & 5.44 & 5.01 \\
H3 robust mean rank & 4.24 & 4.11 \\
H3 fragile mean rank & 7.26 & 6.98 \\
H4 dose-response monotone & yes & yes \\
\midrule
\multicolumn{3}{l}{replicates: H1: yes; H2: yes; H3: yes; H4: yes} \\
\bottomrule
\end{tabular}

\end{table}

\textbf{The preregistered run is void, by its own rule.} One control failed:
Rosenbrock's headroom on the fresh seeds is $0.083$, below the $0.10$ the
preregistration fixed in advance. The document says that a control failure makes
the run void rather than negative, and an exclusion applied after seeing the data
cannot undo that. We report it as void.

We also report what it showed, labelled as the post-hoc sensitivity analysis it
is. Dropping Rosenbrock from both arms, all four claims clear their thresholds:
$\beta(\frag)$ moves from $+0.65$ to $+0.69$ while $\beta(\log_{10}\sigma_e)$
moves from $-0.06$ to $-0.07$ with an interval containing zero; the onset ratio
is $5.01$ against $5.44$; the robust and fragile acquisition families sit at mean
ranks $4.11$ and $6.98$ against $4.24$ and $7.26$; and the dose-response is
monotone in both (Table~\ref{tab:confirmatory}). Two deviations from the
document should be named. The script tested the family contrast over the eight
conditions rather than the twenty benchmarks as preregistered; run as
preregistered it gives $p = 3.6 \times 10^{-5}$ over the twenty benchmarks and
$7.2 \times 10^{-5}$ without Rosenbrock. And the null on the error magnitude
holds under the preregistered additive specification only: with the
magnitude-by-$\optz$ interaction of Section~\ref{sec:mediation} the magnitude
is significant again, so that sub-claim confirms a specification, not a fact
about the error.

One argument is worth making explicitly, because it decides how much the void
should worry a reader. The headroom control exists to stop a benchmark with
nothing to lose from contributing a null that looks like evidence of robustness.
Five of the sub-claims are positive results, and including a low-headroom
benchmark makes those harder to obtain, not easier. The two null-type
sub-claims, that the magnitude's interval contains zero and that Kendall's $W$
is below $0.40$, run the other way, and there the sensitivity estimates are
within $0.01$ of the void run's (interval $[-0.13, 0.09]$ against
$[-0.13, 0.08]$; $W$ $0.272$ against $0.276$). That is a reason to read the
sensitivity analysis as informative; it is not a reason to call the run
confirmed, and we do not.

\FloatBarrier
\section{The landscapes and the descriptor fit}
\label{app:descriptors}

Table~\ref{tab:benchmarks} lists the twenty landscapes with their measured
geometry, and Table~\ref{tab:anchor} places the three human studies of
Section~\ref{sec:cost} on the same axis. Section~\ref{sec:mediation} reports the descriptor model only as a
comparison for $\frag(\sigma_e)$; Table~\ref{tab:descriptors} gives it in full,
with the variance inflation factors that limit what can be read off it.

\begin{table}[htbp]
\centering
\small
\caption{The twenty landscapes and their measured geometry. $\mathrm{opt}_z$ is
how many landscape standard deviations the optimum stands above a random design;
sparsity is the fraction of the box within 10\% of the optimum; ruggedness is
$1$ minus the lag-1 autocorrelation along random walks; tail is
$\sigma/(1.4826\,\mathrm{MAD})$; $\mathrm{frag}(1)$ is the one-shot selection
loss at an error of one landscape standard deviation. Names are the suite's
keys; \texttt{levy\_10} has eleven inputs.}
\label{tab:benchmarks}
\begin{tabular}{lrrrrrrr}
\toprule
landscape & $d$ & $\mathrm{opt}_z$ & sparsity & rugged. & skew & tail & $\mathrm{frag}(1)$ \\
\midrule
\texttt{powell} & 4 & 0.76 & 7.2e-01 & 0.13 & $-$2.33 & 2.52 & 0.18 \\
\texttt{rosenbrock} & 4 & 1.03 & 4.9e-01 & 0.13 & $-$1.06 & 1.13 & 0.18 \\
\texttt{branin} & 2 & 1.25 & 2.4e-01 & 0.13 & $-$1.09 & 1.14 & 0.30 \\
\texttt{steering\_law} & 4 & 1.29 & 7.9e-01 & 0.12 & $-$3.33 & 1.44 & 0.43 \\
\texttt{power\_law\_practice} & 4 & 1.72 & 1.4e-01 & 0.13 & $-$2.02 & 1.72 & 0.63 \\
\texttt{weber\_fechner} & 4 & 2.15 & 5.5e-02 & 0.15 & $-$0.71 & 1.01 & 0.64 \\
\texttt{levy\_10} & 11 & 2.62 & 1.3e-02 & 0.20 & $-$0.54 & 1.00 & 0.80 \\
\texttt{griewank} & 6 & 2.75 & 1.0e-02 & 0.12 & $-$0.26 & 0.98 & 0.85 \\
\texttt{moving\_peaks} & 4 & 2.96 & 1.8e-03 & 0.17 & $-$0.53 & 1.04 & 0.64 \\
\texttt{hartmann\_3} & 3 & 3.05 & 1.6e-02 & 0.14 & +1.14 & 1.49 & 0.50 \\
\texttt{stevens} & 4 & 3.06 & 3.0e-03 & 0.12 & $-$0.17 & 0.97 & 0.75 \\
\texttt{eggholder} & 2 & 3.21 & 6.6e-03 & 0.42 & +0.04 & 1.00 & 0.68 \\
\texttt{yerkes\_dodson} & 4 & 3.54 & 1.9e-03 & 0.12 & +0.27 & 0.98 & 0.73 \\
\texttt{rastrigin} & 4 & 3.64 & 5.8e-04 & 0.58 & $-$0.08 & 0.98 & 0.69 \\
\texttt{schwefel} & 4 & 4.33 & 1.2e-04 & 0.25 & +0.00 & 0.99 & 0.69 \\
\texttt{hicks\_law} & 4 & 5.40 & $<$1.5e-05 & 0.12 & +0.40 & 0.97 & 1.00 \\
\texttt{hartmann\_6} & 6 & 7.97 & 9.2e-05 & 0.14 & +2.67 & 2.74 & 0.53 \\
\texttt{michalewicz} & 5 & 8.05 & $<$1.5e-05 & 0.36 & +1.02 & 0.92 & 0.80 \\
\texttt{ackley} & 4 & 20.10 & $<$1.5e-05 & 0.80 & +3.26 & 1.67 & 0.18 \\
\texttt{shekel} & 4 & 56.77 & $<$1.5e-05 & 0.13 & +3.72 & 1.63 & 0.00 \\
\bottomrule
\end{tabular}

\end{table}


\begin{table}[htbp]
\centering
\caption{Three published human studies placed on the synthetic arm's axis.
$\sigma_f$ is the spread of the study's fitted objective over its design box;
the last column is the measured within-rater noise in those units.}
\label{tab:anchor}
\begin{tabular}{llrrr}
\toprule
study & ratings & $\sigma_f$ & noise / $\sigma_f$ & upper bound \\
\midrule
\texttt{ehmi} & as logged & 0.365 & 0.74 {\scriptsize $[0.48, 0.97]$} & 1.17 \\
\texttt{opticarvis} & as logged, two scales mixed & 0.374 & 3.35 {\scriptsize $[3.03, 3.58]$} & -- \\
\texttt{opticarvis} & one scale & 0.156 & 1.11 {\scriptsize $[0.81, 1.40]$} & 1.34 \\
\texttt{provoice} & as logged, one sign reversed & 0.298 & 1.02 {\scriptsize $[0.20, 1.56]$} (weak) & 2.58 \\
\texttt{provoice} & sign corrected & 0.449 & 0.60 {\scriptsize $[0.00, 0.90]$} (weak) & 2.37 \\
\bottomrule
\end{tabular}

\end{table}

\begin{table}[htbp]
\centering
\caption{Post-onset excess regret on landscape descriptors and the error
magnitude, predictors $z$-scored, cluster-bootstrapped over landscapes with
Benjamini--Hochberg correction. VIF is computed on the design matrix.}
\label{tab:descriptors}
\begin{tabular}{lrcrr}
\toprule
term & $\beta$ & 95\% CI & $p_{\mathrm{FDR}}$ & VIF \\
\midrule
$\log\sigma_e$ & +0.37 & [+0.15, +0.66] & 0.001$^{**}$ & 1.0 \\
$\log\mathrm{opt}_z$ & +0.64 & [+0.24, +0.74] & 0.001$^{**}$ & 9.2 \\
$\log$ tail & +0.20 & [+0.03, +0.33] & 0.008$^{**}$ & 1.2 \\
rugged. & $-$0.09 & [$-$0.38, +0.12] & 0.519 & 1.3 \\
$d$ & $-$0.03 & [$-$0.26, +0.12] & 0.768 & 1.0 \\
\midrule
\multicolumn{5}{l}{$R^2 = 0.526$, 640 cells, 20 landscape clusters} \\
\bottomrule
\end{tabular}

\end{table}

$\optz$ carries a VIF of $9.2$ and regresses on the other descriptors with
$R^2 = 0.89$, so its coefficient and those of sparsity and skew (excluded here,
VIF $5.8$ and $6.3$) are not separately identified. This is why
Appendix~\ref{sec:manipulation} builds landscapes that vary these properties
independently rather than trying to disentangle them by regression. Dimension and ruggedness are null here; the fixed-$\optz$ ladder of
Appendix~\ref{sec:manipulation} shows dimension does raise the cost, so that
null is a limit of the observational design.

\FloatBarrier
\section{Input error: the design that ships}
\label{app:inputerror}

Under an unnoticed slip the optimizer's record holds the proposal $x$, not the
point $x'$ it evaluated, and the design a practitioner deploys is the recorded
one. Table~\ref{tab:inputdeployed} scores both at the end of the run: the best
design the person actually experienced, and the design the log recommends,
scored where the log says it is.

\begin{table}[htbp]
\centering
\caption{Where a slip's cost comes from, at both onsets. \emph{floor}: the
model-free arms, which evaluate displaced points but learn nothing;
\emph{logged}: a learner whose record holds the point actually evaluated;
\emph{unnoticed}: the same learner with the record holding the proposal;
\emph{mislabelling}: the paired difference of the last two, on shared seeds and
slip draws. The early-onset rows also appear in Table~\ref{tab:inputerror}.}
\label{tab:inputmech}
\begin{tabular}{lrrrr}
\toprule
\textsc{slip} & \multicolumn{1}{c}{floor} & \multicolumn{1}{c}{logged} & \multicolumn{1}{c}{unnoticed} & \multicolumn{1}{c}{mislabelling} \\
\midrule
\multicolumn{5}{l}{\textit{error from it.\ 1}} \\
1\% & 0.3 \makebox[3.4em][l]{\scriptsize [0, 1]} & 0.7 \makebox[3.4em][l]{\scriptsize [0, 2]} & 1.5 \makebox[3.4em][l]{\scriptsize [0, 3]} & 0.8 \makebox[3.4em][l]{\scriptsize [0, 2]} \\
5\% & 0.2 \makebox[3.4em][l]{\scriptsize [$-$1, 1]} & 2.8 \makebox[3.4em][l]{\scriptsize [2, 4]} & 5.4 \makebox[3.4em][l]{\scriptsize [3, 8]} & 2.7 \makebox[3.4em][l]{\scriptsize [1, 4]} \\
15\% & 1.2 \makebox[3.4em][l]{\scriptsize [$-$1, 3]} & 7.1 \makebox[3.4em][l]{\scriptsize [5, 10]} & 11.5 \makebox[3.4em][l]{\scriptsize [8, 15]} & 4.4 \makebox[3.4em][l]{\scriptsize [2, 6]} \\
40\% & 6.4 \makebox[3.4em][l]{\scriptsize [0, 13]} & 19.6 \makebox[3.4em][l]{\scriptsize [12, 28]} & 23.2 \makebox[3.4em][l]{\scriptsize [16, 31]} & 3.6 \makebox[3.4em][l]{\scriptsize [2, 6]} \\
\midrule
\multicolumn{5}{l}{\textit{error from it.\ 21}} \\
1\% & 0.0 \makebox[3.4em][l]{\scriptsize [0, 0]} & 0.2 \makebox[3.4em][l]{\scriptsize [0, 1]} & 1.1 \makebox[3.4em][l]{\scriptsize [0, 2]} & 0.9 \makebox[3.4em][l]{\scriptsize [0, 1]} \\
5\% & $-$0.4 \makebox[3.4em][l]{\scriptsize [$-$1, 0]} & 2.0 \makebox[3.4em][l]{\scriptsize [1, 3]} & 3.0 \makebox[3.4em][l]{\scriptsize [2, 5]} & 1.0 \makebox[3.4em][l]{\scriptsize [0, 2]} \\
15\% & $-$0.4 \makebox[3.4em][l]{\scriptsize [$-$1, 1]} & 4.5 \makebox[3.4em][l]{\scriptsize [2, 7]} & 5.0 \makebox[3.4em][l]{\scriptsize [3, 8]} & 0.5 \makebox[3.4em][l]{\scriptsize [0, 1]} \\
40\% & $-$0.6 \makebox[3.4em][l]{\scriptsize [$-$2, 1]} & 6.3 \makebox[3.4em][l]{\scriptsize [3, 10]} & 6.4 \makebox[3.4em][l]{\scriptsize [4, 10]} & 0.1 \makebox[3.4em][l]{\scriptsize [0, 0]} \\
\bottomrule
\end{tabular}

\end{table}

\begin{table}[htbp]
\centering
\caption{Final-iteration excess simple regret under an unnoticed slip, as a
percentage of the achievable improvement, over the model-based acquisitions.
\emph{Evaluated}: the best true value among the points actually touched.
\emph{Deployed}: the true value at the recorded location of the best-rated
trial. Brackets are 95\% cluster-bootstrap intervals over landscapes.}
\label{tab:inputdeployed}
\begin{tabular}{lrrrr}
\toprule
& \multicolumn{2}{c}{from trial 1} & \multicolumn{2}{c}{from trial 21} \\
\cmidrule(lr){2-3} \cmidrule(lr){4-5}
\textsc{slip} & \multicolumn{1}{c}{evaluated} & \multicolumn{1}{c}{deployed} & \multicolumn{1}{c}{evaluated} & \multicolumn{1}{c}{deployed} \\
\midrule
1\% & 1.7 \makebox[3.4em][l]{\scriptsize [0, 3]} & 2.7 \makebox[3.4em][l]{\scriptsize [1, 5]} & 1.5 \makebox[3.4em][l]{\scriptsize [1, 3]} & 2.4 \makebox[3.4em][l]{\scriptsize [1, 4]} \\
5\% & 6.5 \makebox[3.4em][l]{\scriptsize [4, 10]} & 16.4 \makebox[3.4em][l]{\scriptsize [9, 24]} & 3.9 \makebox[3.4em][l]{\scriptsize [2, 6]} & 9.6 \makebox[3.4em][l]{\scriptsize [5, 15]} \\
15\% & 13.4 \makebox[3.4em][l]{\scriptsize [8, 20]} & 34.3 \makebox[3.4em][l]{\scriptsize [24, 45]} & 6.9 \makebox[3.4em][l]{\scriptsize [3, 10]} & 13.4 \makebox[3.4em][l]{\scriptsize [8, 19]} \\
40\% & 24.9 \makebox[3.4em][l]{\scriptsize [16, 34]} & 65.8 \makebox[3.4em][l]{\scriptsize [53, 78]} & 9.0 \makebox[3.4em][l]{\scriptsize [5, 14]} & 13.3 \makebox[3.4em][l]{\scriptsize [8, 19]} \\
\bottomrule
\end{tabular}

\end{table}

\FloatBarrier
\section{Extra trials, by arm}
\label{app:extra}

Table~\ref{tab:extrafull} repeats the count of Section~\ref{sec:cost} for every
arm at the early onset, against a clean ten-trial study within the fifty-trial
budget, with the same one-percent tolerance. Fewer than half the runs are censored in every cell, so every median is exact; the strict count with no tolerance is in \texttt{extra\_runs.csv} beside
each arm. The recorded slip costs fewer trials than the unnoticed one at every magnitude
above $1\%$, and a misclick almost none below $40\%$ of the trials.

\begin{table}[htbp]
\centering
\caption{Median extra trials to match a clean ten-trial study within $1\%$ of
the achievable improvement, error from the first trial, with the share of runs
that never do within 50 trials. Response-error rows are the main sweep over all
ten model-based acquisitions and ten seeds; input-error rows are six
acquisitions and five seeds.}
\label{tab:extrafull}
\begin{tabular}{lrrrr}
\toprule
error & $0.05\sigma$ & $0.25\sigma$ & $1\sigma$ & $5\sigma$ \\
\midrule
\textsc{gaussian} & 0 (1\%) & 1 (8\%) & 9 (26\%) & 26 (44\%) \\
\textsc{bias} & 0 (1\%) & 2 (9\%) & 8 (27\%) & 28 (45\%) \\
\textsc{drift} & 0 (1\%) & 2 (9\%) & 8 (28\%) & 36 (49\%) \\
\textsc{ar}(1) & 0 (1\%) & 1 (4\%) & 3 (22\%) & 23 (42\%) \\
\midrule
error & 1\% & 5\% & 15\% & 40\% \\
\midrule
\textsc{slip} & 0 (3\%) & 3 (7\%) & 8 (15\%) & 24 (40\%) \\
\textsc{slip}, logged & 0 (2\%) & 2 (5\%) & 4 (10\%) & 14 (25\%) \\
\textsc{misclick} & 0 (1\%) & 0 (3\%) & 0 (9\%) & 6 (23\%) \\
\bottomrule
\end{tabular}

\end{table}

\FloatBarrier
\section{Re-scoring the design a study ships}
\label{app:shiprules}

The deployed regret everywhere else in this paper uses one ship rule: the design
with the single best observed rating. To separate search from selection
(Section~\ref{sec:decomp}) we re-score what the same trials would have shipped
under other rules, adding no trial. The loop's own surrogate --- BoTorch's
default single-task GP, normalised over the benchmark's box and standardised ---
is refitted on every logged observation, and each rule picks among the
\emph{visited} designs: the best single rating (the standard), the best mean
rating, the posterior-mean maximiser, the maximiser of the posterior mean less
one or two latent SDs, and, as the ceiling, the visited design with the best true
value. The arm's own logged rule must reproduce the logged final regret to
$10^{-9}$ in every run or the analysis stops; over the $66{,}000$ runs it did, with
no failed fits.

Scored as a share of the deployed cost of \textsc{gaussian} error, the cautious
rule (posterior mean less two latent SDs) recovers $+5.8\%$ $[2, 10]$ at
$1\sigma$ from the first rating, $+12.6\%$ $[4, 23]$ at $1\sigma$ from trial 21,
$+9.7\%$ $[5, 15]$ and $+18.5\%$ $[11, 26]$ at $5\sigma$. It pays a fixed price
of $0.015$ of the achievable improvement whether or not error is present, so
below $0.25\sigma$ it costs $80$--$150\%$ more than the error does and should not
be used without an estimate of the rater's noise. Under \textsc{bias} it survives
only at $1\sigma$ ($+5.5\%$ $[1, 9]$); under \textsc{drift} and \textsc{ar1}
nothing survives correction.

\section{End-of-study procedures}
\label{app:endofstudy}

Both procedures keep $T$ fixed and spend the final trials differently, so each is
replayed exactly from an existing log: truncate at the trial the procedure
starts, refit the loop's GP on that prefix, and simulate only the new trials. The
premise is checked against the re-rating arm, whose earlier trials must equal the
standard run's.

\emph{Tournament($k$)}. Search trials $1..T-k$; the last $k$ are one sitting in
which each of $k$ candidates is shown once. Candidates are the top $k$ visited
designs by the cautious ship rule, or by observed rating; the winner is the best
look, or the posterior mean after conditioning on the looks. Error shared by
every trial of a sitting cancels under comparison, so only the idiosyncratic part
enters: $\sigma$ for \textsc{gaussian}, \textsc{bias} and \textsc{drift}, the
innovation for \textsc{ar1}, and we report the sitting at both that SD and half
of it.

\emph{Confirmation($k$)}. Search trials $1..T-2k$; the last $2k$ rate the chosen
design and the best initial design $k$ times each in counterbalanced order. The
chosen design ships if its \emph{mean} rating over those $k$ is at least the
initial design's, and the study \emph{claims} an improvement only if a one-sided
paired $t$-test at $\alpha=0.05$ prefers it. The two decisions are separate, which
is why this changes little about the design that ships (it recovers $3.2\%$--$3.4\%$
of the cost) and a great deal about what the study may claim.
The standard process claims an improvement in $98.3\%$ of runs and that claim is
false in $15.7\%$ of them, rising to $42.7\%$ at $5\sigma$. Four confirmation
trials take the false-claim rate to $0.5\%$ overall and $1.0\%$ at $5\sigma$, at
the cost of detecting only $50.7\%$ of the real improvements; eight trials give
$0.46\%$ and $65.7\%$.

\section{Deriving the split from the run in hand}
\label{app:budgetrule}

Write $\mu_i, \sigma_i$ for the loop GP's posterior latent mean and SD at the
distinct designs visited by the decision point, ordered by the cautious ship
rule, and $s$ for the idiosyncratic rating SD of a sitting. A $k$-tournament
shows the top $k$ candidates once each, so with $f_i \sim
\mathcal{N}(\mu_i, \sigma_i^2)$ and looks $y_i = f_i + \mathcal{N}(0, s^2)$ the
identification term is $U(k) = \mathbb{E}[f_{\arg\max_i y_i}]$, evaluated by
Monte Carlo. It is not monotone in $k$: a larger candidate set is more likely to
contain the best visited design and more likely to let a worse one win on a lucky
look. The search term prices the forgone trials by the run's own gain in best
rating per trial over the last ten, floored at zero. Every policy --- ship by
posterior mean, ship by the cautious rule, or a tournament of $k$ --- is scored
under the single posterior at $n_0 = T - \max k$, the last trial at which all of
them are still open; scoring each $k$ under its own $T-k$ posterior looks more
faithful and is not, because two GPs fitted on different numbers of points put
their maxima on different scales, and the comparison then picks up that
difference rather than the difference between policies. The sitting's SD is read
from the surrogate's fitted noise, which is the noise of the \emph{standardised}
targets, and multiplied by the outcome transform's scale so that it and the
posterior sit on one axis; without that the rule read a sitting as up to $5.6$
times more precise than it is. The rule reads only ratings and the surrogate; the
true objective enters only when the choice is scored, against the realised regret
of that $(\text{run}, k)$ in the replay.

\section{Budget-neutral arms}
\label{app:budgetneutral}

Every arm below keeps the number of human trials equal to the standard process or
lowers it, and is scored against the standard process by the estimand of
Appendix~\ref{app:adapt}: twenty landscapes, seeds 7--11, both onsets where the
arm has them. Table~\ref{tab:budgetneutral} gives the share of the standard
process's cost of error that each recovers on the design the study deploys.

The pattern is not that everything fails. It is that what fails is the
acquisition, and what works is an explicit model of how the rating was produced.
Clipping the response to the range actually observed recovers $22$--$35\%$ of the
cost of gross faults at a price of exactly zero; a relevance-pursuit likelihood
recovers $26$--$38\%$ when $15\%$ of trials carry a gross fault and loses when
only $5\%$ do; fitting per-rater offsets recovers $32\%$ $[23, 40]$ when raters
alternate and nothing when they hand over in blocks, because a block offset is
confounded with the search's own time trend while an interleaved one is
identifiable; re-anchoring a saturating cap recovers $25\%$ $[17, 32]$; and
removing the factor shared across objectives recovers $25\%$ $[7, 50]$ of the
deployed cost of halo error. A saturating rating scale is itself among the most
expensive faults we measured, costing $185\%$ $[132, 258]$ more than the same
rating noise on a scale that does not saturate.

\begin{table}[htbp]
\centering
\caption{Budget-neutral arms: share of the standard process's cost of error
recovered on the deployed design, with 95\% landscape-bootstrap intervals.
Negative means the change does worse under error than the standard process.
Above the rule, changes to the search; below it, changes to the model of how the
rating was produced.}
\label{tab:budgetneutral}
\small
\begin{tabular}{lr}
\toprule
Arm & Recovered \\
\midrule
Augmented expected improvement & $+2\%$ $[-6, +11]$ \\
Slip-aware acquisition, $0.05$ slip & $-4\%$ $[-15, +4]$ \\
Effort front-loaded to the first ten trials & $-7\%$ $[-20, +4]$ \\
Effort front-loaded to the first twenty & $-8\%$ $[-31, +5]$ \\
Effort on the last ten trials & $-12\%$ $[-24, -2]$ \\
Minimum-distance proposals & $-13\%$ $[-23, -6]$ \\
Slip-aware acquisition, $0.15$ slip & $-18\%$ $[-34, -9]$ \\
Effort split between first and last trials & $-21\%$ $[-33, -10]$ \\
Lower-confidence-bound incumbent & $-27\%$ $[-42, -11]$ \\
Slip-aware acquisition, $0.4$ slip & $-35\%$ $[-54, -16]$ \\
Thompson sampling & $-58\%$ $[-81, -40]$ \\
A lost rating imputed low rather than dropped & $-979\%$ \\
\midrule
Relevance pursuit, $15\%$ of trials spiking at $20$ SD & $+38\%$ $[+24, +51]$ \\
Response clipped to range, $5\%$ spiking at $20$ SD & $+35\%$ $[+23, +46]$ \\
Per-rater offsets, round-robin handover & $+32\%$ $[+23, +40]$ \\
Relevance pursuit, $15\%$ spiking at $5$ SD & $+26\%$ $[+7, +42]$ \\
Re-anchored rating cap & $+25\%$ $[+17, +32]$ \\
Halo backfit across objectives & $+25\%$ $[+7, +50]$ \\
Response clipped to range, $5\%$ spiking at $5$ SD & $+22\%$ $[+10, +35]$ \\
Response clipped to range, $15\%$ spiking at $20$ SD & $+22\%$ $[+4, +38]$ \\
Response clipped to range, $15\%$ spiking at $5$ SD & $+6\%$ $[-7, +18]$ \\
Per-rater offsets, block handover & $+5\%$ $[-5, +14]$ \\
Relevance pursuit, $5\%$ spiking at $20$ SD & $-30\%$ $[-119, +12]$ \\
\bottomrule
\end{tabular}
\end{table}

\section{Comparisons instead of ratings: what the invariance does and does not buy}
\label{app:elicitation}

Preference elicitation is the usual response to an unreliable rater
\citep{chu2005preference, gonzalez2017preferential}, on the argument that a
comparison is invariant to how the rater uses the scale. The argument is exact
and its reach is narrow, and both halves are measurable here. A fault is
\emph{shared} when it acts on every design of one sitting through the same
monotone map $g$; then for strictly increasing $g$, $g(f(a)) > g(f(b))
\Leftrightarrow f(a) > f(b)$, and the comparison is exactly what it would have
been with no fault. A fault drawn afresh for each judgement cancels nothing, and
a comparison inherits the noise of both designs.

We ran two loops at equal human cost on the twenty landscapes, three seeds, $T=20$
judgements: the standard rating loop, and a comparison loop with a Laplace-approximate
\textsc{PairwiseGP} \citep{balandat2020botorch} on which each trial proposes one
design and asks for a comparison with the incumbent, shipping the
posterior-mean maximiser. Ten of the $960$ fits fell back to prior
hyperparameters and are counted, not dropped. Each loop is scored against its own
identically seeded clean twin.

The invariance under a shared strictly monotone fault is not a finding but a
consequence of the construction above, and the run confirms the implementation
rather than the claim: under \textsc{drift} the rating loop's shipped design
differs from its clean twin in $41$ of $60$ runs, at an excess of $0.040$ of the
achievable improvement, and the comparison loop's differs in \emph{none}. What
the run does test is the reach of the invariance, and there it is narrow. Under a
saturating scale, a map that is monotone but
not strictly so, the comparison buys nothing ($-0.3\%$ $[-30, +27]$, an interval
that spans zero): two designs above the cap are indistinguishable, and near the
optimum that is exactly the pair the study must tell apart. Under
\textsc{gaussian} noise, drawn afresh per judgement, the comparison's excess is
$26.5\%$ larger, though the interval $[-79, +15]$ spans zero; the direction is
the one the $\sqrt2$ argument predicts, and the magnitude is not established.
A constant \textsc{bias} present from the first
trial moves neither loop, in no run --- also by construction:
a rating loop that standardises its outputs
and ships an argmax is already shift-equivariant, so that arm tests the rating
loop's invariance rather than the comparison's, and the cost of \textsc{bias}
elsewhere in this paper comes from its arriving \emph{part-way through}.

So changing the elicitation is not a general remedy. It removes, exactly, that
part of the fault which is a strictly monotone reparametrisation of the scale
--- drift, a rater's personal offset, a handover between raters --- and nothing
of the inconsistency, the saturation or the gross faults.

\section{Is the CUSUM the right detector?}
\label{app:detector}

The freeze rule of Section~\ref{sec:budget} detects a mid-study change with a
one-sided CUSUM on the squared one-step-ahead standardised residual of the
loop's own surrogate. A CUSUM is a choice, not a derivation, so we compare it
with two alternatives on exactly the same $13{,}200$ residual streams
($6{,}600$ tuning, $6{,}600$ held out; $6{,}400$ carrying a real mid-study
change). Since the residual is standard normal while the rater behaves and a
rater who becomes noisier inflates its variance, every detector tests
$\mathcal{N}(0,1)$ against $\mathcal{N}(0,\sigma^2)$, $\sigma>1$:
the CUSUM; the generalised likelihood ratio maximised over the unknown change
point, whose statistic $\max_\tau \tfrac{n}{2}(\hat\sigma^2_\tau - 1 -
\log\hat\sigma^2_\tau)$ is the classical answer to that question and what a
CUSUM approximates; and Bayesian online change-point detection
\citep{adams2007bayesian} with a constant hazard and a normal--inverse-gamma
model of the post-change variance, alarming when the posterior probability that
the current run is short exceeds a threshold. Each detector's threshold is tuned
on the tuning seeds to a common false-alarm budget --- an alarm is false when the
rater never changed, or when it fires at or before a late onset --- and scored on
the held-out seeds.

\textbf{The detector is not the bottleneck.} At a $10\%$ budget the GLR detects
$45.5\%$ of changes and the CUSUM $45.4\%$; at $5\%$, $36.8\%$ against $36.9\%$;
at $20\%$, $56.0\%$ against $56.4\%$. The CUSUM matches the classical optimum for
this alternative to within half a percentage point at every budget, so there is
nothing to be gained by replacing it. Bayesian online detection is worse on
detection rate at every budget ($23.0\%$, $32.2\%$, $45.1\%$) but alarms with a
median delay of $2$ trials against the CUSUM's $6$--$9$. Since the freeze
ceiling falls from $82\%$ to $54\%$ after two trials of delay
(Section~\ref{sec:budget}), that trade is not obviously a bad one, but it is a
choice of operating point rather than a better method. What limits the freeze
rule is the signal --- a rater who becomes somewhat noisier midway simply does
not leave much of one --- not the statistic computed from it.

\section{Process adaptations}
\label{app:adapt}

The measurements above say what error costs a fixed process. The practical
question is whether the process can be changed to cost less, and the scoring
has to change with it. Rating a design twice, or re-rating designs before
deploying, also changes the clean run --- a repeated rating of an exact objective
carries no information --- so a process change scored by its excess over its own
clean run is measured against a handicapped reference. We score each change
against the \emph{standard} process instead. Per paired cell (landscape,
acquisition, magnitude, onset, seed), with regret as a fraction of the achievable
improvement,
\begin{equation*}
  \text{recovered} = \frac{R^{\text{std}}_{\text{noisy}} - R^{\text{adapt}}_{\text{noisy}}}
                          {R^{\text{std}}_{\text{noisy}} - R^{\text{std}}_{\text{clean}}},
  \qquad
  \text{price} = R^{\text{adapt}}_{\text{clean}} - R^{\text{std}}_{\text{clean}},
\end{equation*}
aggregated as ratios of landscape means with landscape-bootstrap intervals, on two
responses: the post-onset trajectory regret and the regret of the design the
experimenter would deploy. qKG and replication are different acquisitions, so
there the standard process is the mean over the ten standard acquisitions. Every
arm runs twenty landscapes, seeds 7--11, both onsets and the error process named
in Table~\ref{tab:adapt}, with LogEI and qNEI unless stated; the incumbent and known-noise rows re-score
the arms of Appendix~\ref{sec:noisediag} and~\ref{sec:incumbent} with their own
acquisitions.

\begin{table}[htbp]
\centering
\caption{Share of the standard process's cost of error that each process change
recovers, pooled over magnitudes and onsets, with 95\% landscape-bootstrap
intervals (a star when the interval excludes zero), and the price the change pays
without error in percentage points of the achievable improvement. Negative
recovery means the adapted process does worse under error than the standard
one. The fitted-oracle datasets have no achievable-improvement scale.}
\label{tab:adapt}
\small
\begin{tabular}{lrrrr}
\toprule
& \multicolumn{2}{c}{recovered (\%)} & \multicolumn{2}{c}{price without error} \\
\cmidrule(lr){2-3} \cmidrule(lr){4-5}
adaptation & trajectory & deployed & trajectory & deployed \\
\midrule
\multicolumn{5}{l}{rating error (\textsc{gaussian})} \\
observed-max incumbent & +19$^{*}$ {\scriptsize [10, 27]} & $-$10$^{*}$ {\scriptsize [$-$15, $-$5]} & +0.1 & +0.2 \\
GP given the true noise variance & $-$3 {\scriptsize [$-$16, 8]} & $-$3 {\scriptsize [$-$11, 5]} & +0.4 & $-$0.0 \\
first ten proposals rated twice & $-$88$^{*}$ {\scriptsize [$-$113, $-$67]} & $-$14$^{*}$ {\scriptsize [$-$21, $-$8]} & +5.1 & +1.6 \\
\quad + observed-max incumbent (LogEI) & $-$93$^{*}$ {\scriptsize [$-$127, $-$58]} & $-$44$^{*}$ {\scriptsize [$-$57, $-$32]} & +4.7 & +1.1 \\
last six trials re-rate the top three & $-$1$^{*}$ {\scriptsize [$-$2, $-$1]} & +4 {\scriptsize [0.04, 9]} & +0.1 & +0.8 \\
knowledge gradient & $-$50$^{*}$ {\scriptsize [$-$83, $-$21]} & $-$39$^{*}$ {\scriptsize [$-$61, $-$18]} & +4.5 & +5.3 \\
re-rate the incumbent every second trial & $-$114$^{*}$ {\scriptsize [$-$155, $-$76]} & $-$12 {\scriptsize [$-$31, 5]} & +7.1 & +6.0 \\
\multicolumn{5}{l}{unnoticed \textsc{slip}} \\
noisy-input GP & $-$28$^{*}$ {\scriptsize [$-$43, $-$13]} & $-$19$^{*}$ {\scriptsize [$-$31, $-$9]} & $-$0.1 & $-$1.5 \\
last six trials re-rate the top three & 0 {\scriptsize [$-$0.3, 0.2]} & $-$6 {\scriptsize [$-$15, 1]} & +0.1 & +0.8 \\
\multicolumn{5}{l}{\textsc{misclick}} \\
Student-$t$ surrogate & $-$2 {\scriptsize [$-$38, 31]} & $-$13 {\scriptsize [$-$50, 16]} & +0.2 & +1.5 \\
\multicolumn{5}{l}{changing what the rater is asked or when} \\
the rater reports their precision & +16 {\scriptsize [$-$4, 36]} & +15$^{*}$ {\scriptsize [6, 23]} & +0.2 & $-$0.1 \\
each proposal judged beside the best & $-$18$^{*}$ {\scriptsize [$-$25, $-$12]} & $-$27$^{*}$ {\scriptsize [$-$36, $-$19]} & +0.0 & +0.0 \\
every fifth trial rates an anchor & $-$6 {\scriptsize [$-$23, 14]} & +1 {\scriptsize [$-$8, 10]} & +0.8 & +0.7 \\
the first five proposals rated late & $-$13 {\scriptsize [$-$44, 14]} & +1 {\scriptsize [$-$11, 11]} & +1.6 & +0.2 \\
a ship-rule acquisition & $-$287$^{*}$ {\scriptsize [$-$366, $-$219]} & $-$107$^{*}$ {\scriptsize [$-$159, $-$64]} & +21.8 & +24.4 \\
\multicolumn{5}{l}{fitted oracles, rating error} \\
first ten proposals rated twice & $-$161 {\scriptsize [$-$3504, $-$49]} & $-$8 {\scriptsize [$-$17, 6]} & -- & -- \\
\bottomrule
\end{tabular}

\end{table}

\textbf{Two of the paper's own ablations, re-scored.} The observed-max incumbent
and the known-noise GP change the clean run too, so their shares in
Appendix~\ref{sec:noisediag} and~\ref{sec:incumbent} are also excess over their
own clean runs. Against the standard process the conclusions stand, with one
addition. The incumbent recovers $19\%$ $[10, 27]$ of the trajectory cost at a
price of $0.1$ points without error, most of it from the first rating ($+10\%$
to $+67\%$ across magnitudes), and cuts the extra trials needed to match a clean
standard ten-trial study at $1\sigma$ from $8$ ($21\%$ never) to $6$ ($13\%$),
and at $5\sigma$ from $36$ ($48\%$) to $21$ ($39\%$). The design it deploys is
worse, though, by $10\%$ of the cost $[5, 15]$, mostly at the late onset, for a
reason we have not isolated. Giving the GP the true noise variance recovers
nothing on either response ($-3\%$).

\textbf{Rating early designs twice makes things worse.} The onset effect says the
first trials are where error is expensive, so replicating them is the natural
fix. It is not one: rating each of the first ten proposals twice recovers
$-88\%$ of the trajectory cost and $-14\%$ of the deployed design's
(Table~\ref{tab:adapt}), and it needs more extra trials to match a clean standard
ten-trial study, a median of $8$ at $1\sigma$ against $6$ for the standard
process. In a fifty-trial study the ten designs it gives up cost more than
averaging two ratings, which only divides the error variance by two, buys; the GP
already pools neighbouring observations. Adding the observed-max incumbent of
Appendix~\ref{sec:incumbent} does not rescue it ($-44\%$ on the deployed design),
and the same replication on the three fitted-oracle datasets points the same way
($-8\%$ on the deployed design, interval $[-17, 6]$). Re-rating the incumbent every second trial fails for the same reason ($-61\%$ on
the trajectory). qKG does not replicate: it helps with error from the first trial
and hurts with late error and on the deployed design, for a reason we have not
tested (Appendix~\ref{sec:designchoices}).

\textbf{Re-evaluating before deploying helps, a little.} Spending the last six
trials re-rating the three designs with the best mean rating, and deploying by
mean, costs almost nothing in the trajectory ($-1\%$) and recovers $4\%$ of the
deployed design's cost, $[0.04, 9]$, reaching $18\%$ $[5, 29]$ when large error
arrives late. Under an unnoticed slip it recovers nothing ($-6\%$, $[-15, 1]$):
the re-ratings land on displaced points too, so averaging them does not correct
the record.

\textbf{A noisy-input GP helps small early slips and hurts large late ones.}
Inflating each point's observation variance by the squared gradient of the
posterior mean times the slip variance recovers $76\%$ $[27, 245]$ of the cost of
a $1\%$ slip from the first trial and $17\%$ of a $5\%$ slip, but at the late onset it recovers $-129\%$ and $-205\%$ of the cost of $15\%$ and
$40\%$ slips. A likely reason, which we have not tested, is that by then the search
sits near a steep peak whose flanks have the largest gradients, so the
observations that locate the peak are discounted most.
Pooled it recovers $-28\%$ of the trajectory cost.

\textbf{A heavy-tailed likelihood does not absorb misclicks.} A misclick records
the rating of a random design against the proposed one, an outlier rather than
noise, so a Student-$t$ likelihood should discount it. A variational Student-$t$ GP
in place of the standard one recovers $-2\%$ $[-38, 31]$ of the trajectory cost
and $-13\%$ $[-50, 16]$ of the deployed design's, at a price of $0.2$ and $1.5$
points without error. Where misclicks cost most, on $40\%$ of trials from the
first, it recovers $-1\%$ $[-10, 6]$ of the trajectory cost, and the median run
still needs $8$ extra trials to match a clean standard ten-trial study, as it does
without it ($24\%$ never, against $28\%$). Misclicks cost little in the first
place, so there is little for any model to recover.

\FloatBarrier
\section{The one-shot loss from a pilot}
\label{app:pilot}

Equation~\ref{eq:frag} is defined on the exact objective at 256 quasi-random
designs, which a practitioner does not have. For each landscape and each of the
ten seeds we fit the study's own GP to the first $k$ trials of a clean LogEI run
and evaluate Equation~\ref{eq:frag} on its posterior mean, with the same 256
designs and the same $4{,}000$ error draws as the exact value
(Table~\ref{tab:pilot}).

\begin{table}[htbp]
\centering
\caption{Spearman correlations across the twenty landscapes of $\frag$ computed from
the GP posterior mean after $k$ clean trials (averaged over ten seeds) with the
exact $\frag$ and with the measured cost at the early onset. Pooled rows use all
eighty landscape-by-magnitude cells.}
\label{tab:pilot}
\begin{tabular}{lrrrr}
\toprule
& \multicolumn{1}{c}{with exact frag} & \multicolumn{3}{c}{with measured cost} \\
\cmidrule(lr){2-2} \cmidrule(lr){3-5}
frag computed from & pooled & pooled & $1\sigma$ & $5\sigma$ \\
\midrule
first 10 trials & 0.94 & 0.85 & 0.32 & 0.61 \\
first 20 trials & 0.96 & 0.85 & 0.20 & 0.66 \\
first 50 trials & 0.96 & 0.85 & 0.15 & 0.63 \\
exact objective & 1.00 & 0.84 & 0.42 & 0.45 \\
\bottomrule
\end{tabular}

\end{table}

Ten trials are enough. Pooled over magnitudes, $\frag$ from the first ten trials
correlates $0.94$ with the exact value and $0.85$ with the measured cost, against
$0.84$ for the exact $\frag$ itself; twenty and fifty trials add almost nothing.
Within a single magnitude the ranking of landscapes is weak for the exact $\frag$
too ($0.08$ to $0.45$), so the predictive power is mostly across magnitudes; at
$5\sigma$ the pilot underestimates $\frag$ by about a third on average (a bias of $-0.78$)
but ranks
landscapes by cost better than the exact version ($0.61$ against $0.45$).

\FloatBarrier
\section{Does noise ever help?}
\label{app:beneficial}

A noisy observation can in principle help, by pushing an over-exploiting
optimizer out of a local basin. At one seed per cell, $25$ of the $640$
landscape-by-condition cells had negative mean excess regret with $d_z < -0.2$,
noise apparently improving the outcome, which is the kind of finding that
invites a mechanism.

At ten seeds it is one cell: \texttt{michalewicz} under \textsc{ar}(1) error at
$0.05\sigma$ from iteration $21$, mean excess $-0.136$, $d_z = -0.21$ over $100$
paired runs. One cell in $640$ is fewer than the fifteen or so that a null family of $640$
tests would flag at this threshold. We report it as a null and note the
sampling-variability lesson: the twenty-four cells that disappeared between one
seed and ten were all in the direction that would have made the better story.

\end{document}
